\ifx\danslelivre\undefined
\RequirePackage{luatex85}
\documentclass[10pt]{../fichiers-configuration-latex/smfart}
\usepackage{polyglossia}
%\usepackage[francais]{babel}
\usepackage{fontspec}
\usepackage[math-style=ISO,bold-style=upright,sans-style=upright]{unicode-math}
%\usepackage[mathletters]{ucs}
%\usepackage[utf8]{inputenc}
%\usepackage[T1]{fontenc}
\usepackage[unicode,colorlinks=true,linkcolor=blue,citecolor=Sepia,urlcolor=magenta]{hyperref}
\setmainlanguage{french}
\setmainfont[Ligatures=TeX]{Linux Libertine O}
% variante possible : TeXGyrePagella ?
\setmathfont{XITS Math}

% maths
\usepackage{amsmath}
%\usepackage{amssymb}
%\usepackage{mathrsfs}
%\usepackage{stmaryrd}
%\usepackage{euscript}

% couleurs
\usepackage[usenames,dvipsnames]{xcolor}

% graphiques
\usepackage{graphicx}
\usepackage[all]{xy}
% Xy-pic 3.8.9's PDF driver assumes that \fontdimen8\textfont3 is a
% legacy TeX rule width.  Under LuaTeX with unicode-math/XITS Math it is
% instead 6.62pt, which turns every arrow shaft into a thick black bar.
% Override only xypdf's private cached width; do not alter the math font.
\makeatletter
\AtBeginDocument{%
  \@ifundefined{xP@preclw}{}{%
    \edef\xP@preclw{0.4pt}%
    \edef\xP@lw{\xP@dim\xP@preclw}%
    \PackageInfo{sga4-xypdf-compat}{Xy-pic PDF line width set to \xP@preclw}%
  }%
}
\makeatother
%\usepackage{tikz}
%\usetikzlibrary{matrix,arrows,calc}
%\tikzset{column sep=3em,row sep=2em,text height=1.5ex, text depth=0.25ex,>=angle 90}
%\tikzset{description/.style={fill=white,inner sep=1pt}}

% hyperliens, 'etiquettes
\usepackage{xr-hyper}
\usepackage{../fichiers-configuration-latex/smfthm}
%\ifx\cacheretiquettes\undefined
%\usepackage[inline]{showlabels} \else \fi % pour rendre visibles les 'etiquettes
%\usepackage{../fichiers-configuration-latex/makeidx} % pour faire un index
%\usepackage[pagebackref,unicode]{hyperref}  % pour voir o`u sont cit'ees les r'ef'erences
\hypersetup{unicode,colorlinks=true,linkcolor=blue,citecolor=Sepia,urlcolor=magenta,
  hypertexnames=false}
%\hypersetup{breaklinks=true} %permet le retour `a la ligne dans les liens trop longs. Incompatible avec pdflatex
%\renewcommand*{\backref}[1]{{\small [cit\'e page(s)~#1.]}}

\usepackage{xspace}
\usepackage{pifont} % pour expose XVII

\renewcommand{\thefootnote}{\roman{footnote}} % pour mettre les notes en bas de page sous la forme (ii) et non (2), qui ressemble trop (?) a un exposant.

\newcounter{pageoriginal}
\newcommand{\pageoriginale}{\refstepcounter{pageoriginal}\marginpar{\footnotesize\textbf{\thepageoriginal}}\xspace}
\let\pageoriginaled\pageoriginale

\usepackage{ifthen}
\newboolean{original}
\def\sisi{\ifthenelse{\boolean{original}}}
\setboolean{original}{false}

\newcounter{cnote}
\def\nde#1{\ifthenelse{\boolean{original}}{}{\stepcounter{cnote}\def\thefootnote{\arabic{cnote}}\protect\footnote[\value{cnote}]{N.D.E.: #1}}}
\def\ndemark{\ifthenelse{\boolean{original}}{}{\stepcounter{cnote}\def\thefootnote{\arabic{cnote}}\protect\footnotemark[\value{cnote}]}}
\def\ndetext#1{\ifthenelse{\boolean{original}}{}{\def\thefootnote{\arabic{cnote}}\protect\footnotetext[\value{cnote}]{N.D.E.: #1}}}

\setmainlanguage{french}
\setotherlanguage{english}

% unicode-math deliberately disables the legacy \Arrowvert glyph.  The
% source uses it once as a rotated equality marker; \Vert has the intended
% appearance and keeps the historical diagram buildable with LuaLaTeX.
\renewcommand{\Arrowvert}{\Vert}


%%
%% Pour afficher le numéro de version Git (exige que LuaTeX soit appelé avec --shell-escape)
%%

\catcode`\%=12
\directlua{
function command_output(cmd)
  local output_pipe = io.popen(cmd)
  local line_iter = output_pipe and output_pipe:lines()
  local line = line_iter and line_iter()
  if output_pipe then output_pipe:close() end
  return line
end
local git_command1 = [[git log --pretty=format:"%h %ai" -1]]
local git_command2 = [[if git diff --quiet HEAD ; then echo clean ; else echo dirty ; fi]]
function git_status()
  local status = command_output(git_command1)
  if status then
    tex.tprint({luatexbase.catcodetables.CatcodeTableLaTeX, string.char(0x5c).."texttt{"}, {-2,status}, {luatexbase.catcodetables.CatcodeTableLaTeX, "}"})
    if not (command_output(git_command2) == "clean") then tex.sprint(-2, " (modifications non validées)") end
  else
    tex.sprint(-2, "?")
  end
end
}
\catcode`\%=14
\def\showgitstatus{\directlua{git_status()}}


% macros pour LaTeX
\renewcommand{\mathsf}{}

% Police différente pour U+2113 SCRIPT SMALL L
\setmathfont[range="2113]{latinmodern-math.otf}

% Give every explicit SGA label its own destination and preserve the printed
% number supplied as the optional argument.  The former definition inherited
% whichever counter anchor happened to precede it and ignored #1 entirely.
\makeatletter
\def\etiquette[#1]#2{%
  \phantomsection
  \protected@edef\@currentlabel{#1}%
  \label{#2}%
}
\makeatother
\def\eqhyperref[#1]#2{(\hyperref[#1]{#2})}
\def\pagehyperref[#1]#2{(page \hyperref[#1]{#2})}
\def\hyperrefext[#1][#2]#3{\hyperref[#2]{#1 #3}}
\def\eqhyperrefext[#1][#2]#3{\hyperref[#2]{#1 #3}}     %cf. exposes 17
\def\neweqhyperrefext[#1][#2]#3{\hyperref[#2]{#1 #3}}  %cf. exposes 17

\def\shat#1{\protect\widehat{#1}}
\def\stilde#1{\protect\widetilde{#1}}
%%%%%%%%%%%%%%%%%%


\def\etcn{etc.\xspace}
\def\etco{etc...\xspace}
\def\etc{\sisi{\etco}{\etcn}}
\def\ldotsn{\sisi{...}{\etcn}}

\def\cf{{cf.}\kern.3em}
\def\Cf{{Cf.}\kern.3em}
\def\ie{{i.e.}\kern.3em}
\def\iev{{i.e.},\xspace}
\def\loccito{loc.\kern2pt cit.\xspace}
\def\loccitn{\textit{loc.\kern2pt cit.}\xspace}
\def\loccit{\sisi{\loccito}{\loccitn}}
\def\EGA{\textup{EGA}\kern.2em}
\def\SGA{\textup{SGA}\kern.3em}

\def\quoi{}

\def\quav{\quad,\quad}
\def\quaet{\quad \text{et}\quad}
\def\fibreproduct#1#2#3{#1 ×_{#3} #2}
\let\fprod\fibreproduct
\def\ccheck#1#2{\mbox{${\check{\check{\rm #1}}}^{\lower2pt\hbox{$\scriptstyle #2$}}$}}
\def\sheaf#1{\protect\mathcal{#1}}
\def\vieuxsheaf#1{{#1}} % cf. problème avec sheaf{a\Phi_i} par exemple...
\DeclareMathOperator{\Hom}{Hom}
\DeclareMathOperator{\can}{can}
\DeclareMathOperator{\Ind}{Ind}
\DeclareMathOperator{\ind}{ind}
%\DeclareMathOperator{\Hhomtop}{\mathcal{H}\!\mathit{omtop}}
\DeclareMathOperator{\Hhomtop}{\mathbf{Homtop}}
\DeclareMathOperator{\Point}{Point}
\DeclareMathOperator{\Pro}{Pro}
\DeclareMathOperator{\Top}{Top}
\DeclareMathOperator{\Mod}{Mod}
\DeclareMathOperator{\Cart}{Cart}
%\DeclareMathOperator{\Mmorsite}{\mathcal{M}\!\mathit{orsite}}
\DeclareMathOperator{\Mmorsite}{\mathbf{Morsite}}
\DeclareMathOperator{\loc}{loc}
\DeclareMathOperator{\Loc}{Loc}
\DeclareMathOperator{\Pic}{Pic}
\DeclareMathOperator{\Et}{Et}
\DeclareMathOperator{\card}{card}
\DeclareMathOperator{\Ens}{Ens}
\DeclareMathOperator{\h}{\underline{H}}
\DeclareMathOperator{\Fonct}{Fonct}
\DeclareMathOperator{\Ker}{Ker}
\DeclareMathOperator{\but}{but}
\DeclareMathOperator{\Ob}{Ob}
\DeclareMathOperator{\ob}{ob}
\DeclareMathOperator{\Lib}{Lib}
\DeclareMathOperator{\esj}{esj}
\DeclareMathOperator{\id}{id}
\DeclareMathOperator{\Id}{Id}
\DeclareMathOperator{\ad}{ad}
\DeclareMathOperator{\Iim}{Im}
\DeclareMathOperator{\pr}{pr}
\DeclareMathOperator{\Cat}{Cat}
\DeclareMathOperator{\car}{car}
\DeclareMathOperator{\SysInd}{SysInd}
\DeclareMathOperator{\Sup}{Sup}
\DeclareMathOperator{\Coker}{Coker}
\DeclareMathOperator{\Filt}{Filt}
\DeclareMathOperator{\acc}{acc}
\DeclareMathOperator{\sur}{sur}
\DeclareMathOperator{\cat}{cat}
\DeclareMathOperator{\Ord}{Ord}
\DeclareMathOperator{\Gr}{Gr}
\DeclareMathOperator{\Cov}{Cov}
\DeclareMathOperator{\Caf}{Caf}
\DeclareMathOperator{\coker}{coker}
\DeclareMathOperator{\coim}{coim}
\DeclareMathOperator{\im}{im}
\DeclareMathOperator{\Ouv}{Ouv}
\DeclareMathOperator{\diag}{diag}
\DeclareMathOperator{\Esp}{Esp}
\DeclareMathOperator{\TOP}{TOP}
\DeclareMathOperator{\Unif}{Unif}
\DeclareMathOperator{\Iind}{\mathbf{Ind}}
\DeclareMathOperator{\Hhomcart}{\mathbf{Homcart}}
\DeclareMathOperator{\Hhomtopess}{\mathbf{Homtopess}}
\DeclareMathOperator{\Ppoints}{\mathbf{Points}}
\DeclareMathOperator{\Ppoint}{\mathbf{Point}}
\DeclareMathOperator{\Ttors}{\mathbf{Tors}}
\DeclareMathOperator{\Ssex}{\mathbf{Sex}}
\DeclareMathOperator{\Ffib}{\mathbf{Fib}}
\DeclareMathOperator{\Vvois}{\mathbf{Vois}}
\DeclareMathOperator{\Ppointess}{\mathbf{Pointess}}
\DeclareMathOperator{\scrHom}{\mathbf{Hom}}
\DeclareMathOperator{\sheafPic}{\mathbf{Pic}}
\DeclareMathOperator{\sheafAut}{\mathbf{Aut}}
\DeclareMathOperator{\Hhomtoplocan}{\mathbf{Homtoplocan}}
\DeclareMathOperator{\Hhomtopan}{\mathbf{Homtopan}}
\DeclareMathOperator{\Pptgeom}{\mathbf{Ptgéom}}
\DeclareMathOperator{\Morsite}{Morsite}
\DeclareMathOperator{\Fib}{Fib}
\DeclareMathOperator{\Fais}{Fais}
\DeclareMathOperator{\kar}{kar}
\DeclareMathOperator{\Groupes}{Groupes}
\DeclareMathOperator{\Restr}{Restr}
\DeclareMathOperator{\Isom}{Isom}
\DeclareMathOperator{\Prol}{Prol}
\DeclareMathOperator{\Res}{Res}
\DeclareMathOperator{\Tors}{Tors}
\DeclareMathOperator{\esp}{esp}
\DeclareMathOperator{\es}{es}
\DeclareMathOperator{\od}{od}
\DeclareMathOperator{\gr}{gr}
\DeclareMathOperator{\sob}{sob}
\DeclareMathOperator{\dfn}{dfn}
\DeclareMathOperator{\Sch}{Sch}
\DeclareMathOperator{\ouv}{ouv}
\DeclareMathOperator{\cons}{cons}
\DeclareMathOperator{\Nerf}{Nerf}
\DeclareMathOperator{\Points}{Points}
\DeclareMathOperator{\Espsob}{Espsob}
\DeclareMathOperator{\Homtop}{Homtop}
\DeclareMathOperator{\Pt}{Pt}
\DeclareMathOperator{\Quot}{Quot}
\DeclareMathOperator{\glob}{glob}
\DeclareMathOperator{\Homfib}{Homfib}
\DeclareMathOperator{\Qouv}{Qouv}
\DeclareMathOperator{\Quotouv}{Quotouv}
\DeclareMathOperator{\ens}{ens}
\DeclareMathOperator{\ab}{ab}
\DeclareMathOperator{\an}{an}
\DeclareMathOperator{\fs}{fs}
\DeclareMathOperator{\red}{red}
\DeclareMathOperator{\Ab}{Ab}
\DeclareMathOperator{\cosq}{cosq}
\DeclareMathOperator{\Ann}{Ann}
\DeclareMathOperator{\Zar}{Zar}


%%%%%%%%%%%%
%%%%% SGA4-2
\DeclareMathOperator{\Sex}{Sex}
\DeclareMathOperator{\vois}{vois}
\DeclareMathOperator{\Vois}{Vois}
\DeclareMathOperator{\Ext}{Ext}
\DeclareMathOperator{\ext}{ext}
\DeclareMathOperator{\Fer}{Fer}
\DeclareMathOperator{\Add}{Add}
\DeclareMathOperator{\Cocart}{Cocart}
\DeclareMathOperator{\cosk}{cosk}
\DeclareMathOperator{\plat}{plat}

\DeclareMathOperator{\faisc}{faisc}
\DeclareMathOperator{\sk}{sk}
\DeclareMathOperator{\ch}{ch}

\def\HH{\mathbf{H}}
\def\qc{\mathrm{qc}}
\def\et{\mathrm{et}}
\def\coh{\mathrm{coh}}

\def\Hvr#1{%
\begin{array}[b]{l}
{\kern2pt\scriptstyle r}\\[-8pt]
{\kern1pt\scriptstyle \vee}\\[-5pt]
H^{#1}
\end{array}}

\def\Hvv#1{%
\begin{array}[b]{l}
\kern2pt\rotatebox{90}{${\scriptscriptstyle\ll}$}\\[-6pt]
H^{#1}
\end{array}}

\def\Hvi#1{%
\begin{array}[b]{l}
{\kern0.2pt\scriptscriptstyle \infty}\\[-8pt]
{\kern1pt\scriptstyle \vee}\\[-5pt]
H^{#1}
\end{array}}

\def\Hv#1{%
\begin{array}[b]{l}
{\kern1pt\scriptstyle \vee}\\[-5pt]
H^{#1}
\end{array}}

\DeclareMathOperator{\cart}{cart}
\DeclareMathOperator{\Ffaisc}{\mathbf{Faisc}}


\def\Ev#1{%
\begin{array}[b]{l}
{\kern1pt\scriptstyle \vee}\\[-5pt]
E'_{#1}
\end{array}}

\DeclareMathOperator{\Cont}{Cont}
\DeclareMathOperator{\Ccont}{\mathscr{C}\!\mathit{ont}}
\DeclareMathOperator{\Site}{Site}
\newcommand\MM{\mathbf{M}}
\def\limtop#1{{\displaystyle{\varprojlim_{#1}\!\mbox{\rm top }}}}
\def\Limtop{\mathop{\varprojlim\!\mbox{\rm top}}}
\DeclareMathOperator{\spec}{spec}
\def\ss{\scriptstyle{\mathscr{S}}}
\DeclareMathOperator{\Gal}{Gal}
\def\smallm{\scriptstyle{\mathscr{M}}}
\DeclareMathOperator{\End}{End}
\DeclareMathOperator{\ass}{ass}
\DeclareMathOperator{\zar}{zar}
\DeclareMathOperator{\Ogg}{Ogg}
\DeclareMathOperator{\res}{res}
\DeclareMathOperator{\tr}{tr}
\DeclareMathOperator{\Supp}{Supp}
\DeclareMathOperator{\codim}{codim}
\DeclareMathOperator{\rk}{rk}
\DeclareMathOperator{\cur}{cur}
\DeclareMathOperator{\cdqc}{cdqc}
\newcommand\DD{\mathbf{D}}
\DeclareMathOperator{\cl}{cl}
\def\JJ{\mathbf{J}}
\def\FF{\mathbf{F}}
\def\BB{\mathbf{B}}
\DeclareMathOperator{\Alget}{Alget}
\DeclareMathOperator{\Ass}{Ass}
\DeclareMathOperator{\supp}{supp}
\DeclareMathOperator{\rad}{rad}
\DeclareMathOperator{\prof}{prof}
\DeclareMathOperator{\cint}{int}
\DeclareMathOperator{\sic}{sic}
\DeclareMathOperator{\Tor}{Tor}
\DeclareMathOperator{\tors}{tors}
\DeclareMathOperator{\Sym}{Sym}
\DeclareMathOperator{\parf}{parf}
\DeclareMathOperator{\tor}{tor}
\DeclareMathOperator{\defn}{defn}
\DeclareMathOperator{\Div}{Div}
\DeclareMathOperator{\Pol}{Pol}
\DeclareMathOperator{\Fex}{Fex}
\DeclareMathOperator{\Lie}{Lie}
\DeclareMathOperator{\tot}{tot}
\DeclareMathOperator{\Aut}{Aut}
\DeclareMathOperator{\HOM}{HOM}
\DeclareMathOperator{\pch}{pch}
\DeclareMathOperator{\Ch}{Ch}
\DeclareMathOperator{\EXT}{EXT}
\DeclareMathOperator{\PIC}{PIC}
\DeclareMathOperator{\Rf}{Rf}
\DeclareMathOperator{\inj}{inj}
\newcommand\VV{\mathbf{V}}

\def\Xv{%
\begin{array}[b]{l}
{\kern1pt\scriptstyle \vee}\\[-5pt]
X
\end{array}}

%%%%%%%
\newcommand\TT{\mathbf{T}}
\newcommand\XX{\mathbf{X}}
\newcommand\EE{\mathbf{E}}
%%%%%%

\def\qedsymbol{\textup{C.Q.F.D.}}
\let\CQFD\qed
\let\cqfd\qed

%%%%%\bigdot
\font\bigsymb=cmsy10 at 4pt
\def\bigdot{{\kern1.2pt\raise 1.5pt\hbox{\bigsymb\char15}}}
%%%%%

\def\Lright{\mathop{\vcenter{\hbox{\hbox{$\longrightarrow$}\lower8pt \hbox{\kern-16pt\hbox{$\longrightarrow$}}}}}}

\def\Aa{{\mathcal A}}
\def\Bb{{\mathcal B}}
\def\Cc{{\mathcal C}}
\def\Dd{{\mathcal D}}
\def\Ee{{\mathcal E}}
\def\Ff{{\mathcal F}}
\def\Gg{{\mathcal G}}
\def\Hh{{\mathcal H}}
\def\Ii{{\mathcal I}}
\def\Jj{{\mathcal J}}
\def\Kk{{\mathcal K}}
\def\Ll{{\mathcal L}}
\def\Mm{{\mathcal M}}
\def\Nn{{\mathcal N}}
\def\Oo{{\mathcal O}}
\def\Pp{{\mathcal P}}
\def\Qq{{\mathcal Q}}
\def\Rr{{\mathcal R}}
\def\Ss{{\mathcal S}}
\def\Tt{{\mathcal T}}
\def\Uu{{\mathcal U}}
\def\Vv{{\mathcal V}}
\def\Ww{{\mathcal W}}
\def\Xx{{\mathcal X}}
\def\Yy{{\mathcal Y}}
\def\Zz{{\mathcal Z}}

\newcommand{\CC}{\mathbf{C}}
\newcommand{\GG}{\mathbf{G}}
\newcommand{\KK}{\mathbf{K}}
\newcommand{\LL}{\mathbf{L}}
\newcommand{\NN}{\mathbf{N}}
\newcommand{\PP}{\mathbf{P}}
\newcommand\QQ{\mathbf{Q}}
\newcommand\RR{\mathbf{R}}
\newcommand\ZZ{\mathbf{Z}}

\def\to{\mathchoice{\longrightarrow}{\rightarrow}{\rightarrow}{\rightarrow}}
\def\from{\mathchoice{\longleftarrow}{\leftarrow}{\leftarrow}{\leftarrow}}
\def\mto{\mathchoice{\longmapsto}{\mapsto}{\mapsto}{\mapsto}}
\def\mfrom{\mathrel{\from\joinrel\mapstochar\,}}
\def\hto{\mathrel{\lhook\joinrel\to}}
\def\hfrom{\mathrel{\from\joinrel\rhook}}

\def\implies{\mathchoice{\Longrightarrow}{\Rightarrow}{\Rightarrow}{\Rightarrow}}
\def\implied{\mathchoice{\Longleftarrow}{\Leftarrow}{\Leftarrow}{\Leftarrow}}
\def\iff{\mathchoice{\Longleftrightarrow}{\Leftrightarrow}{\Leftrightarrow}{\Leftrightarrow}}

\def\To#1{\mathchoice{\xrightarrow{\labelstyle~~#1~}}{\stackrel{#1}{\longrightarrow}}{}{}}
\def\Mto#1{\mathrel{\mapstochar\To{#1}}}
\def\From#1{\mathchoice{\xleftarrow{\labelstyle~~#1~}}{\stackrel{#1}{\longleftarrow}}{}{}}

\newcommand{\isomto}{\overset{\sim}{\longrightarrow}}
\newcommand{\isomfrom}{\overset{\sim}{\longleftarrow}}

\newcommand{\IFF}{\unskip~$\iff$~\ignorespaces}
\newcommand{\THEN}{\unskip~$\implies$~\ignorespaces}

\let\leq\leqslant
\let\geq\geqslant
\let\le\leqslant
\let\ge\geqslant

\def\cf{{cf.}\kern.3em}
\def\Cf{{Cf.}\kern.3em}
\def\ie{{i.e.}\kern.3em}
\def\iev{{i.e.},\xspace}
%\def\loccito{loc.\kern2pt cit.\xspace}
%\def\loccitn{\textit{loc.\kern2pt cit.}\xspace}
%\def\loccit{\sisi{\loccito}{\loccitn}}
\def\EGA{\textup{EGA}\kern.2em}
\def\SGA{\textup{SGA}\kern.3em}

\let\oldS\S
\def\S{\oldS\kern .15em}
\def\ptbl{.\kern .2em}
\def\resp{resp.\kern.3em}
\def\Exp{Exp.\kern.3em}

\let\bar\overline
\let\hat\widehat
\let\tilde\widetilde
\let\dpl\displaystyle
\let\ldots\dots
\let\emptyset\varnothing
\let\skpt\skippointrait
\let\epsilon\varepsilon

\def\sfrac#1#2{{#1}/{#2}}

\DeclareMathOperator{\cd}{cd}
%\DeclareMathOperator{\Hhom}{\mathcal{H}\!\mathit{om}}
\DeclareMathOperator{\Hhom}{\mathbf{Hom}}
%\DeclareMathOperator{\Eext}{\mathcal{E}\!\mathit{xt}}
\DeclareMathOperator{\Eext}{\mathbf{Ext}}
\DeclareMathOperator{\Spec}{Spec}
\DeclareMathOperator{\Tr}{Tr}

\let\emptyset=\varnothing

%%%%%%%%%%%%%%%
%%% expose XVII
%%%%%%%%%%%%%%%

\makeatletter
\let\old@optexpose\optexpose
\newcommand{\empty@optexpose}[1]{}
\newcommand{\paren@optexpose}[1]{\old@optexpose{#1}(}
\newcommand{\refnoexpose}[1]{\let\optexpose\empty@optexpose\ref{#1}\let\optexpose\old@optexpose}%\refnoexpose{X.cor5.2} engendre "5.2" au lieu de "X 5.2"
\newcommand{\neweqref}[1]{\let\optexpose\paren@optexpose\textup{\ref{#1})}\let\optexpose\old@optexpose}%\neweqref{XII.eq4.1} engendre "XII (4.1)" au lieu de "(XII 4.1)"
\makeatother
\newcommand{\refvb}[2][VI${}_{\text{B}}$ ]{\sisi{#1\refnoexpose{#2}}{\ref{#2}}}

\newcommand{\leftexp}[2]{{\vphantom{#2}}^{#1}{#2}}
\newcommand{\AAA}{\mathbf{A}}
\newcommand{\Sf}{\mathfrak{S}}
\DeclareMathOperator*{\quotprojlim}{\text{``$\varprojlim$''}}
\DeclareMathOperator*{\quotinjlim}{\text{``$\varinjlim$''}}
\DeclareMathOperator{\Fl}{F\sisi{\ell}{l}}
\DeclareMathOperator{\Pl}{P\sisi{\ell}{l}}
%\DeclareMathOperator{\cl}{c\sisi{\ell}{l}}
\DeclareMathOperator{\tdf}{tdf}
\DeclareMathOperator{\rg}{rg}
\newcommand{\iInd}{\sisi{I}{i}nd}
\newcommand{\iind}{ind}
\newcommand{\ppro}{pro}
\newcommand{\tTor}{\sisi{T}{t}or}
\newcommand{\ttor}{tor}
\newcommand{\chap}{Chap.}
\newcommand{\Ffex}{\Ff\mathit{ex}}
\newcommand{\sheafk}{\underline{k}}
\newcommand{\sheafn}{\underline{n}}
\newcommand{\sheafGamma}{\underline{\Gamma}}

\sisi{
\DeclareMathOperator{\simarrow}{\rule{.15mm}{1.8mm}\kern -3.4pt\rightsquigarrow}
\newcommand{\sheafell}{\underline{\ell}}
\newcommand{\sheafZZ}{\underline{\ZZ}}
}{
\newcommand{\nderef}[1]{la N.D.E.~(\ref{#1}) page~\pageref{#1}}
}

%%%%%%%%%%%%%%%%%
%%%%% expose XVII
%%%%%%%%%%%%%%%%%



\externaldocument{../01/01.tex}
\externaldocument{../02/02.tex}
\externaldocument{../03/03.tex}
\externaldocument{../04/04.tex}
\externaldocument{../05/05.tex}
\externaldocument{../05bis/05bis.tex}
\externaldocument{../06/06.tex}
\externaldocument{../07/07.tex}
\externaldocument{../08/08.tex}
\externaldocument{../09/09.tex}
\externaldocument{../10/10.tex}
\externaldocument{../11/11.tex}
\externaldocument{../12/12.tex}
\externaldocument{../13/13.tex}
\externaldocument{../14/14.tex}
\externaldocument{../15/15.tex}
%\externaldocument{../16/16.tex}
\externaldocument{../17/17.tex}
\externaldocument{../18/18.tex}
\externaldocument{../19/19.tex}

\setboolean{original}{false}

\begin{document}

\begin{center}
XVI. Théorème de changement de base par un morphisme lisse, et applications\\
M. Artin 
\end{center}

\selectlanguage{english}version : \showgitstatus\selectlanguage{french}
\tableofcontents

\else
\chapter{Théorème de changement de base par un morphisme lisse, et applications}
%\addtocontents{toc}{M. Artin}
\begin{center}
M. Artin 
\end{center}
\fi

\setcounter{pageoriginal}{206}

\section{Le théorème de changement de base par un morphisme
  lisse}\etiquette[1]{XVI.sec1}\pageoriginale

\begin{enonce*}{Théorème 1.1}\etiquette[1.1]{XVI.thm1.1}
Soient $\LL\subset \PP$, $n\in \NN$, et
$$\vcenter{
\xymatrix@C=1.5cm@R=1cm{
X   \ar[d]^{f}& X' \ar[l]_{g'} \ar[d]^{f'}\\
Y     & \ar[l]_{g} Y'
}}\leqno{(*)}
$$
un diagramme cartésien, avec $f$ quasi-compact et quasi-séparé\nde{=cohérent XXX}, et $g$
universellement localement $0$-acyclique (\resp
\sisi{univ. loc.}{universellement localement} $1$-asphérique pour $\LL$, \resp
\sisi{univ. loc.}{universellement localement} $n$-acyclique pour $\LL$) \emph{(\hyperrefext[XV][XV.def1.11]{1.11}).} Alors
pour chaque faisceau $F$ d'ensembles (\resp de
\emph{ind}-$\LL$-groupes, \resp abélien de
$\LL$-torsion) sur $X$, le morphisme de changement de base
\emph{(XII 4.1.2)}
$$
\varphi^{q}:g^{\ast}(R^{q}f_{\ast})F\longrightarrow
(R^{q}f'_{\ast}){g'}^{\ast}F
$$
est bijectif pour $q=0$ (\resp $q\leq 1$, \resp $q\leq n$).
\end{enonce*}

En particulier, en appliquant \hyperrefext[XV][XV.thm2.1]{2.1}, on trouve le

\begin{enonce*}{Corollaire 1.2}\etiquette[1.2]{XVI.cor1.2}
(Théorème de changement de base par un morphisme lisse). Supposons que
  le morphisme $g$ de $(\ast)$ soit lisse, que $f$ soit quasi-compact
  et quasi-séparé, et que $\LL\subset \PP$ soit
  l'ensemble complémentaire à l'ensemble des caractéristiques
  résiduelles de $Y$. Alors pour chaque faisceau $F$ d'ensembles
  (\resp de \emph{ind}-$\LL$-groupes, \resp abélien de
  $\LL$-torsion) sur $X$, le morphisme de
  changement\pageoriginale de base $\varphi^{q}$ ci-dessus est
  bijectif pour $q=0$ (\resp pour $q=0$, $1$, \resp pour chaque $q$).
\end{enonce*}

\noindent
\emph{Démonstration de} \hyperref[XVI.thm1.1]{1.1}. Traitons d'abord le cas où
$f:X\rightarrow Y$ est quasi-projectif: ce n'est qu'une conjonction du
théorème de changement de base pour un morphisme propre \hyperrefext[XII][XII.thm5.1]{5.1} et de
la définition \hyperrefext[XV][XV.prop1.9]{1.9}. En effet, on a un diagramme commutatif
$$
\xymatrix@=1.25cm{
 X \ar[rr]^{i} \ar[rd]_{f}&        & \bar{X}\ar[ld]^{\bar{f}}\\
                         &   Y    &
}
$$
où $\bar{f}$ est projectif, donc propre, et où $i$ est une immersion
ouverte. Or on sait \hyperrefext[XII][XII.thm5.1]{5.1} que le théorème de changement de base est
vrai pour $\bar{f}$ quel que soit $g:Y'\rightarrow Y$. Puisque
$\fprod{g}{\bar{Y}}{Y}:\fprod{Y'}{\bar{Y}}{Y}\rightarrow \bar{Y}$ est
encore universellement $0$-acyclique (\resp\ldots), on se ramène
immédiatement par (\hyperrefext[XII][XII.prop4.4]{4.4} (ii)) à démontrer le théorème pour le
morphisme $i$, c'est-à-dire, on est ramené au cas où $f$ est une
immersion ouverte. Alors le théorème est conséquence de la définition
\hyperrefext[XV][XV.prop1.10]{1.10} (ii).

Traitons maintenant le cas général. L'assertion est locale sur $Y$, et
on peut donc supposer $Y$ affine.

\begin{enonce*}{Lemme 1.3}\etiquette[1.3]{XVI.lem1.3}
On peut de plus supposer $X$ affine.
\end{enonce*}

\noindent
\emph{Démonstration}. Procédant comme dans \hyperrefext[XII][XII.lem6.1]{6.1}, on est réduit,
pour prouver \hyperref[XVI.thm1.1]{1.1} pour un $f$ donné, à le prouver dans la situation
déduite\pageoriginale de la situation donnée par des changements de
base $\widetilde{Y}\rightarrow Y$, $\widetilde{Y}'\rightarrow Y'$,
pour un morphisme local $\widetilde{g}:\widetilde{Y}'\rightarrow
\widetilde{Y}$ de localisés stricts de $Y'$, $Y$ induit par $g$, et à
prouver que les homomorphismes
\begin{equation*}
H^{q}(\widetilde{X},F)\longrightarrow H^{q}(\widetilde{X}',F')\quad
(\mbox{\rm où } F'={g'}^{\ast}(F))\tag{$\ast$}
\end{equation*}
sont bijectifs pour $q\leq n$, où $\widetilde{X}=\fprod{X}{Y}{Y}$,
$\widetilde{X}'=\fprod{X'}{Y'}{Y'}$. Cela nous ramène donc au cas où
$Y$, $Y'$, $g$ sont strictement locaux, et à prouver dans ce cas la
bijectivité des applications précédentes, admettant que
\hyperref[XVI.thm1.1]{1.1} est prouvé 
pour $f$ affine. Or soit $(X_i)$ un recouvrement fini de $X$ par des
ouverts affines $X_i$, et soit $Z$ le schéma somme des $X_i$, qui est
donc affine et muni d'un morphisme surjectif $h:Z\rightarrow X$, d'où
un morphisme $h':Z'\rightarrow X'$. Notons que $h$ est séparé, et
qu'il est affine lorsque $f$ est séparé. Ceci dit, appliquant le lemme
de descente \hyperrefext[XII][XII.prop6.8]{6.8}, on voit que pour prouver la bijectivité des
applications $(\ast)$ pour $q\leq n$ et pour tout $F$, il suffit de
prouver la bijectivité des applications correspondantes
$H^{q}(Z,h^{\ast}(F))\rightarrow H^{q}(Z',{h'}^{\ast}(F'))$, à
condition que les homomorphismes de changement de base, pour
$Z\rightarrow X$ et le changement de base $X'\rightarrow X$, soient
des isomorphismes en dimension $\leq n$. Or si $f$ est séparé, donc
$h$ est affine, il en est ainsi par hypothèse, ce qui prouve \hyperref[XVI.thm1.1]{1.1}
lorsque $f$ est supposé séparé. Dans le cas général, on peut alors
appliquer le résultat précédent à $h$ qui est toujours séparé, et on
conclut encore que \hyperref[XVI.thm1.1]{1.1} est vrai pour $f$, ce qui prouve \hyperref[XVI.lem1.3]{1.3}.

\begin{enonce*}[remark]{Remarque}
On pourrait aussi invoquer la suite spectrale de Leray pour le
recouvrement ouvert $(X_i)$ de $X$, et ses variantes non
commutatives,\pageoriginale mais cette méthode, qui n'est pas
essentiellement différente du recours à \hyperrefext[XII][XII.prop6.8]{6.8}, a le désavantage de
nous obliger à distinguer à nouveau les trois cas habituels, travail
qui a déjà été fait dans \loccit
\end{enonce*}

Pour achever la démonstration du théorème dans le cas général, il
suffit de ramener la démonstration au cas où $f:X\rightarrow Y$ est
affine et de type fini, donc quasi-projectif. D'après \hyperref[XVI.lem1.3]{1.3}, on peut
supposer $X$ affine. Soit alors $X=\varprojlim X_{\alpha}$, où
$f_{\alpha}:X_{\alpha}\rightarrow Y$ sont des schémas affines et de
type fini sur $Y$. On termine avec la technique de passage à la limite
habituelle (\hyperrefext[VII][VII.sec5]{5}).

\begin{enonce*}[remark]{Compléments 1.4}\etiquette[1.4]{XVI.compl1.4}
\begin{enumerate}
\renewcommand{\theenumi}{\alph{enumi}}
\renewcommand{\labelenumi}{\theenumi)}
\item Nous laissons au lecteur le soin de constater que la même
  démonstration donne encore une conclusion lorsqu'on suppose que $g$
  est universellement localement $(-1)$-acyclique (\hyperrefext[XV][XV.rems1.14]{1.14}): dans ce
  cas, l'homomorphisme de changement de base
$$
\phi^{0}:g^{\ast}(f_{\ast}(F))\longrightarrow f'_{\ast}({g'}^{\ast}(F))
$$
est \emph{injectif} pour tout faisceau d'ensembles $F$ sur
$X$. D'autre part, supposons à nouveau $n\geq 0$, on peut compléter
\hyperref[XVI.thm1.1]{1.1} par l'énoncé d'injectivité suivant, contenu également dans la
démonstration qui précède: si $g$ est universellement localement
$0$-acyclique, alors pour tout faisceau en groupes $F$ sur $X$,
l'homomorphisme de changement de base $\phi^{1}$ de \hyperref[XVI.thm1.1]{1.1} est un
monomorphisme; si $g$ est universellement localement $n$-acyclique
pour $\LL$, alors l'homomorphisme $\phi^{n+1}$ de \hyperref[XVI.thm1.1]{1.1} est un
monomorphisme pour tout faisceau abélien de $\LL$-torsion $F$
sur $X$.

\item Il\pageoriginale est sans doute possible de donner également une
  conclusion d'injectivité analogue à partir de l'hypothèse de
  $1$-asphéricité locale pour $\LL$ de $g$, en introduisant les
  invariants de $2$-cohomologie non commutative de la thèse de \textsc{Giraud},
  comparer \hyperrefext[XII][XII.subsec5.11]{5.11} \footnote[1]{C'est ce qui est effectivement
  établi dans le livre de J. Giraud.}; la même question se pose
  également pour (\hyperref[XVI.cor1.6]{1.6}) et (\hyperref[XVI.cor2.3]{2.3}) ci-dessous.
\end{enumerate}
\end{enonce*}

%raghu
\begin{enonce*}{Corollaire 1.5}\etiquette[1.5]{XVI.cor1.5}
Soit $K/k$ une extension de corps, et soit
$\LL=\PP-\{p\}$, où $p= \car k$. Alors $\Spec
(K)\rightarrow \Spec (k)$ est universellement localement acyclique
pour $\LL$ et universellement localement $1$-asphérique pour
$\LL$; si $k$ et $K$ sont séparablement clos, alors le
morphisme précédent est également universellement acyclique pour
$\LL$ et universellement $1$-asphérique pour $\LL$.
\end{enonce*}

Notons tout de suite la conséquence suivante de (\hyperref[XVI.cor1.5]{1.5}):

\begin{enonce*}{Corollaire 1.6}\etiquette[1.6]{XVI.cor1.6}
Soient $K/k$ une extension de corps séparablement clos, $X$ un schéma,
$F$ un faisceau d'\sisi{ensemble}{ensembles} (\resp de ind-$\LL$-groupes, \resp
de groupes abéliens de $\LL$-torsion) sur $X$, où
$\LL=\PP-\{p\}$, $p=\car k$. Alors l'application
$$
H^{q}(X,F)\longrightarrow H^{q}(X_{K},F_{K})
$$
est bijective pour $q=0$ (\resp pour $q\leq 1$, \resp pour tout $q$).
\end{enonce*}

Cela résulte en effet aussitôt du fait que $\Spec(K)\rightarrow
\Spec(k)$ est universellement acyclique pour $\LL$ et
universellement $1$-asphérique pour $\LL$, et des
définitions.\sisi{---~}{~}Notons qu'on peut interpréter aussi l'homomorphisme
précédent sur les $H^{q}$ comme s'identifiant, par \hyperrefext[VIII][VIII.cor2.4]{2.4}\sisi{}{,}
à l'homomorphisme de changement de base
$$
g^{\ast}(R^{q}f_{\ast}(F))\longrightarrow
R^{q}{f'}_{\ast}({g'}^{\ast}(F)){\quoi},
$$
où\pageoriginale $f:X\rightarrow \Spec(k)$, $f':X_K\rightarrow
\Spec(K)$, $g:\Spec(K)\rightarrow \Spec(k)$ et $g':X_K\rightarrow X$
sont les morphismes évidents. Donc lorsque $X$ est quasi-compact et
quasi-séparé, la bijectivité de ces homomorphismes peut aussi être
considérée, grâce à \hyperref[XVI.thm1.1]{1.1}, comme conséquence du fait que
$\Spec(K)\rightarrow \Spec(k)$ est universellement \emph{localement}
acyclique pour $\LL$ et universellement localement
$1$-asphérique pour $\LL$. Ce dernier fait implique donc déjà
\hyperref[XVI.cor1.6]{1.6} pour $X$ quasi-compact et quasi-séparé, ce qui suffit
manifestement à entraîner que $\Spec (K)\rightarrow \Spec(k)$ est
universellement acyclique pour $\LL$ et universellement
$1$-asphérique pour $\LL$ (\hyperrefext[XV][XV.def1.7]{1.7}, \hyperrefext[XV][XV.prop1.6]{1.6} (i)).

Ceci montre donc que pour prouver \hyperref[XVI.cor1.5]{1.5}, il suffit de prouver la
première assertion de \hyperref[XVI.cor1.5]{1.5}. Quitte à passer à la clôture parfaite de
$k$, ce qui est licite par \hyperrefext[VIII][VIII.thm1.1]{1.1},
on peut alors supposer $k$ \emph{parfait}. Toute extension $K$ de $k$ est limite inductive de ses
sous-algèbres $A$ de type fini sur $k$, et $k$ étant parfait, quitte à
localiser un tel $A$, on peut le supposer lisse, de sorte que $K$
apparaît comme limite inductive filtrante de sous-algèbres lisses. Ces
dernières étant universellement localement acycliques pour
$\LL$ et universellement localement $1$-asphériques pour
$\LL$ en vertu de \hyperrefext[XV][XV.thm2.1]{2.1}, on conclut grâce à \hyperrefext[XV][XV.prop1.13]{1.13} (ii),
\cqfd

\begin{enonce*}[remark]{Remarque 1.7}\etiquette[1.7]{XVI.rem1.7}
On comparera le théorème d'invariance \hyperref[XVI.cor1.5]{1.5} à
\hyperrefext[XII][XII.cor5.4]{5.4}, où on n'a pas eu à supposer le faisceau de torsion
envisagé premier aux caractéristiques résiduelles, mais où en revanche
on doit supposer $X$ \emph{propre} sur $k$. Déjà dans le cas de
$H^{1}$ $(X,\ZZ/p\ZZ)$, où $X$\pageoriginale est une courbe algébrique
affine (par exemple la droite affine) sur un corps algébriquement clos
$k$ de car. $p>0$, l'analogue des énoncés précédents devient faux, à
cause des phénomènes de « ramification immodérée\nde{=sauvage XXX} » à l'infini,
impliquant que dans un tel cas la classification des revêtements
étales principaux de groupe $\ZZ/p\ZZ$ est essentiellement « continue ».
\end{enonce*}

\section{Théorème de spécialisation des groupes de cohomologie}\etiquette[2]{XVI.sec2}

\begin{enonce*}{Théorème 2.1}\etiquette[2.1]{XVI.thm2.1}
Soient $\LL\subset \PP$, $n\in \NN$, et soit
$f:X\rightarrow Y$ un morphisme de présentation finie propre
\footnote[1]{Pour une hypothèse moins restrictive que celle de
  propreté permettant d'obtenir les mêmes conclusions, \cf \SGA 5 II 3
et \SGA 1 XIII.} et localement $0$-acyclique (\resp
\sisi{loc.}{localement} $1$-asphérique pour $\LL$, \resp \sisi{loc.}{localement} $n$-acyclique pour
$\LL$). Soit $F$ un faisceau d'ensembles (\resp de
\emph{ind}-$\LL$-groupes, \resp abélien de
$\LL$-torsion) constructible et localement constant sur
$X$. Alors les $R^{q}f_{\ast}F$ sont constructibles et \emph{localement
constants} pour $q=0$ (\resp pour $q=0$, $1$, \resp pour $q\leq n$) et
pour tout point géométrique $\bar{y}$ de $Y$, on a
$$
(R^{q}f_{\ast}F)_{\bar{y}}\simeq
H^{q}(X_{\bar{y}},F_{\bar{y}}){\quoi},
$$
pour ces mêmes valeurs de $q$. 
\end{enonce*}

\begin{enonce*}[remark]{Remarque}
Notons que la dernière assertion n'est que le théorème de changement
de base pour un morphisme propre \hyperrefext[XII][XII.cor5.2]{5.2}.
\end{enonce*}

Compte tenu de \hyperrefext[XV][XV.thm2.1]{2.1}, on a immédiatement le corollaire suivant:\pageoriginale

\begin{enonce*}{Corollaire 2.2}\etiquette[2.2]{XVI.cor2.2}
(Théorème de spécialisation pour les groupes de cohomologie). Soit
  $f:X\rightarrow Y$ un morphisme propre et lisse, et
  $\LL\subset \PP$ l'ensemble complémentaire à
  l'ensemble des caractéristiques résiduelles de $Y$. Soit $F$ un
  faisceau d'ensembles (\resp de ind-$\LL$-groupes, \resp
  abélien de $\LL$-torsion) constructible sur $X$. Alors les
  $R^{q}f_{\ast}(F)$ sont constructibles et \emph{localement
  constants} pour $q=0$ (\resp pour $q=0$, $1$, \resp pour tout $q$),
  et pour tout point géométrique $\bar{y}$ de $Y$, on a
$$
R^{q}f_{\ast}(F)_{\bar{y}}\simeq H^{q}(X_{\bar{y}},F_{\bar{y}})
$$
pour ces mêmes valeurs de $q$.
\end{enonce*}

\begin{enonce*}[remark]{Remarque}
Prenant $n=1$, et traduisant en termes de groupes fondamentaux, on
retrouve (\SGA 1 X 3.8) comparant les groupes fondamentaux des fibres
d'un morphisme propre et lisse.
\end{enonce*}

\noindent
\emph{Démonstration de} \hyperref[XVI.thm2.1]{2.1}. Puisqu'on sait déjà que les
$R^{q}f_{\ast}(F)$ sont constructibles \hyperrefext[XIV][XIV.thm1.1]{1.1}, il suffit grâce à 
\hyperrefext[IX][IX.prop2.11]{2.11} de prouver que pour tout morphisme de spécialisation
$\bar{y}_1\rightarrow \bar{y}_{0}$ de points géométriques de $Y$,
le morphisme de spécialisation correspondant (\hyperrefext[VIII][VIII.subsec7.7]{7.7})
\begin{equation*}
R^{q}f_{\ast}(F)_{\bar{y}_{0}}\longrightarrow
R^{q}f_{\ast}(F)_{\bar{y}_1}\tag{$\ast$}
\end{equation*}
est bijectif pour les valeurs de $q$ envisagées. Or on a à ce sujet le
résultat un peu plus général suivant (où il est inutile de
supposer\pageoriginale $f$ de présentation finie et $F$ constructible):

\begin{enonce*}{Corollaire 2.3}\etiquette[2.3]{XVI.cor2.3}
Soient $\LL\subset \PP$, $n\in \NN$,
$f:X\rightarrow Y$ un morphisme propre et localement $0$-acyclique
(\resp localement $1$-asphérique pour $\LL$, \resp localement
$1$-acyclique pour $\LL$), $F$ un faisceau d'ensembles (\resp
de ind-$\LL$-groupes, \resp abélien de $\LL$-torsion)
sur $X$, qui est localement constant. Alors pour tout morphisme de
spécialisation $\bar{y}_1\rightarrow \bar{y}_0$ de points géométriques
de $Y$, l'homomorphisme de spécialisation correspondant $(\ast)$
ci-dessus est bijectif pour $q=0$ (\resp pour $q=0$, $1$, \resp pour
$q\leq n$). De plus, dans le cas $n=0$ \ie $g$ localement
$0$-acyclique, pour tout faisceau de groupes localement constant $F$
sur $X$, l'homomorphisme $R^{1}f_{\ast}(F)_{\bar{y}_0}\rightarrow
R^{1}f_{\ast}(F)_{\bar{y}_1}$ est injectif; si $n$ est quelconque,
pour tout faisceau abélien de $\LL$-torsion $F$ localement
constant sur $X$, l'homomorphisme de spécialisation $(\ast)$ est
injectif pour $q=n+1$.
\end{enonce*}

Il existe un schéma strictement local intègre $Y'$, et un morphisme
$g:Y'\rightarrow Y$ appliquant le point fermé $y'_0$ en $y_{0}$, le point
générique en $y'_1$: il suffit par exemple de prendre d'abord le
localisé strict de $Y$ relativement à $\bar{y}_0$, puis son
sous-schéma fermé intègre défini par un point au-dessus de $y_1$. De
plus, quitte à remplacer $Y'$ par son normalisé dans une clôture
algébrique de son corps des fonctions, on peut supposer $Y$ normal et
$k(y'_1)$ algébriquement clos, de sorte que $y'_0$ et $y'_1$ sont des
points géométriques de $Y'$. Utilisant le théorème de changement de
base (\hyperrefext[XII][XII.thm5.1]{5.1}) pour $(f,F,g)$, on est ramené, pour prouver \hyperref[XVI.cor2.3]{2.3}, à le
faire en remplaçant\pageoriginale $Y$, $\bar{y}_0$, $\bar{y}_1$ par
$Y'$, $y'_0$, $y'_1$, ce qui nous ramène à prouver 2.3 dans le cas où
$Y$ est normal et strictement local, et où $\bar{y}_0$ et $\bar{y}_1$
sont respectivement son point fermé $y_0$ et son point générique
$y_1$. D'ailleurs, utilisant \hyperrefext[XII][XII.sec4]{4}, on voit que l'homomorphisme de
spécialisation $(\ast)$ s'identifie alors à l'homomorphisme
$$
H^{q}(X,F)\longrightarrow H^{q}(X_1,g^{\ast}_1(F)){\quoi},
$$
où $X_1=X_{y_1}$ et où $g:y_1\rightarrow Y$ et $g_1:X_1\rightarrow X$
sont les morphismes canoniques. Notre assertion provient alors du
lemme suivant (où on ne suppose plus $f$ propre):

\begin{enonce*}{Lemme 2.4}\etiquette[2.4]{XVI.lem2.4}
Avec les notations préliminaires de $2.3$, abandonnons l'hypothèse de
propreté sur $f$, supposons en revanche $Y$ normal intègre, de point
générique $y_1$ tel que $k(y_1)$ soit séparablement clos, et soient
$X_1=X_{y_1}$, $g_1:X_1\rightarrow X$ le morphisme canonique. Alors
l'homomorphisme
\begin{equation*}
H^{q}(X,F)\longrightarrow H^{q}(X_1,g^{\ast}_1(F))\tag{${\ast}{\ast}$}
\end{equation*}
est bijectif pour $q=0$ (\resp bijectif pour $q=0$, $1$, \resp
bijectif pour $q=n$, injectif pour $q=n+1$). De plus, dans le cas
$n=0$ \ie $g$ localement acyclique, pour tout faisceau en groupes $F$
sur $X$, l'homomorphisme $({\ast}{\ast})$ est injectif pour $q=1$.
\end{enonce*}

Un argument immédiat (\hyperrefext[XV][XV.prop1.6]{1.6} (i) $\Rightarrow$ (ii)) montre qu'il
suffit de prouver la relation
$$
F\xrightarrow{\sim} g_{1_{\ast}}g_1{}^{\ast}(F){\quoi},
$$
et\pageoriginale pour $n\geq 1$, les relations supplémentaires
$$
R^{q}g_{1_{\ast}}g_1{}^{\ast}(F)=0\quad\mbox{\rm pour}\quad 1\leq
q\leq n{\quoi}.
$$
Ces relations étant locales sur $X$ pour la topologie étale, et $F$
étant localement constant pour la topologie étale, on peut supposer
$F$ constant, donc de la forme $G_X$, où $G$ est un ensemble (\resp un
ind-$\LL$-groupe, \resp un groupe abélien de
$\LL$-torsion). Notons d'autre part que l'hypothèse que \sisi{$X$}{$Y$} est
normal implique $g_{\ast}(G_{y_1})\xleftarrow{\sim}G_Y$ (\hyperrefext[IX][IX.lem2.14.1]{2.14.1}),
et l'hypothèse que $k(y_1)$ est séparablement clos implique que les
$R^{q}g_{\ast}(G_{y_1})$ sont nuls pour $q\geq 1$. D'autre part,
l'hypothèse que $g$ est localement $0$-acyclique (\resp\ldots)
implique que la formation des $R^{q}g_{\ast}(G_y)$ commute au
changement de base $f:X\rightarrow Y$ pour $q\leq n$, compte tenu du
fait que l'on peut supposer $Y$ affine, et qu'alors $y_1$ apparaît
comme limite projective de ses voisinages ouverts affines, qui sont
quasi-compacts dans $Y$, de sorte qu'on peut appliquer la définition
\hyperrefext[XV][XV.def1.11]{1.11} et la théorie du passage à la limite \hyperrefext[VII][VII.sec5]{5}.
On en conclut qu'on a bien
$$
G_X\xrightarrow{\sim} g_{1_{\ast}}(G_{X_1})\quav
R^{q}g_{1_{\ast}}(G_{X_{1}})=0\quad\mbox{\rm pour}\quad 1\leq q\leq
n{\quoi},
$$
ce qui achève la démonstration de \hyperref[XVI.lem2.4]{2.4} et par suite de \hyperref[XVI.thm2.1]{2.1}.

\begin{enonce*}[remark]{Remarque}
La démonstration qui précède, via \hyperref[XVI.lem2.4]{2.4}, plus simple que notre
démonstration initiale, est due à M. Lubkin.
\end{enonce*}

\begin{enonce*}{Corollaire 2.5}\etiquette[2.5]{XVI.cor2.5}
Sous les conditions de $2.1$, supposons $Y$ connexe, et soit $q\leq
n$. Alors les $H^{q}(X_{\bar{y}},F)$, pour les points
géométriques\pageoriginale $y$ de $Y$, sont tous isomorphes entre eux.
\end{enonce*}

\section{Le théorème de pureté cohomologique relatif}\etiquette[3]{XVI.sec3}

\begin{enonce*}{Définition 3.1}\etiquette[3.1]{XVI.def3.1}
Soit $S$ un schéma. On appelle \emph{$S$-couple lisse} $(Y,X)$ une
$S$-immersion fermée $i:Y\rightarrow X$ de $S$-préschémas lisses,
c'est-à-dire, un diagramme commutatif de préschémas
$$
\xymatrix@=1.25cm{
Y \ar[rr]^{i} \ar[rd]_{h}&     &X\ar[ld]^{f}\\
  & S   &
}
$$
tel que $f$ et $h$ soient lisses et que $i$ soit une immersion
fermée. On notera $U=X-Y$, $j:U\rightarrow X$ l'immersion ouverte, et
$g:U\rightarrow S$ le morphisme structural. On appelle
\emph{codimension} de $(Y,X)$ en un point $y\in Y$ la codimension en
$y$ de la fibre $Y_s$ dans $X_s$, où $s=h(y)$.
\end{enonce*}

La codimension est donc une fonction localement constante sur
$Y$. Nous indiquerons par la phrase « $(Y,X)$ est de
codimension~$c$ » que cette fonction est même constante, \sisi{à}{de} valeur~$c$. % ?

Soient $(Y,X)$ un $S$-couple lisse et $y\in Y$. Rappelons (\SGA II
4.10) qu'il existe des entiers $m$ et $n$, un voisinage $X'$ de $y$
dans $X$, et un morphisme étale
$$
\phi:X'\longrightarrow \EE^n_S=\Spec \Oo_S[t_1,\ldots,t_n]
$$
tels\pageoriginale que $Y'=Y\cap X'$ soit l'image inverse du
sous-schéma fermé $\EE^m_S$ défini par les équations
$$
t_{m+1}=\ldots=t_n=0{\quoi}.
$$
On peut exprimer ce résultat en disant que \emph{localement pour la
  topologie étale, chaque $S$-couple lisse $(Y,X)$ est isomorphe au
  couple standard $(\EE^m_S, \EE^n_S)$ pour $m$ et $n$
  convenables,} où
$$
\EE^r_S=\Spec \Oo_S[t_1,\ldots,t_r]{\quoi}.
$$

Dans le présent numéro, on se propose de calculer les faisceaux de
cohomologie locale $\HH^q_Y(X,F)=(R^{q}i^{!})F$ (\sisi{\hyperrefext[V][V.sec4]{4}}{\hyperrefext[V][V.sec6]{6}} et \hyperrefext[VIII][VIII.subsec6.6]{6.6})
pour un $S$-couple lisse $(Y,X)$, à valeurs dans un faisceau abélien de
torsion localement constant $F$ premier aux caractéristiques
résiduelles. Nous commençons par des considérations préliminaires:

\begin{enonce*}{Proposition 3.2}\etiquette[3.2]{XVI.prop3.2}
Soit
$$
\xymatrix@=1.25cm{
U \ar[rr]^{j} \ar[rd]_{g}&     &X\ar[ld]^{f}\\
  & S   &
}
$$
un diagramme commutatif, où $f$ est lisse, et $j$ est une immersion
ouverte telle que la fibre $U_s$ soit dense dans $X_s$ pour chaque
$s\in S$. Soit $F$ un faisceau d'ensembles sur $S$. Alors le morphisme
canonique
$$
f^{\ast}F\longrightarrow j_{\ast}g^{\ast}F=j_{\ast}j^{\ast}f^{\ast}F
$$
est bijectif.
\end{enonce*}

\noindent
\emph{Démonstration}.\pageoriginale On se réduit facilement au cas où $U$ est
rétrocompact dans $X$. Soit $\bar{s}$ un point géométrique de $S$,
$\bar{x}$ un point géométrique de $X$ au-dessus de $S$,
$\widetilde{f}:\widetilde{X}\rightarrow \widetilde{S}$ le morphisme
des localisés stricts de $S$, $X$ en des points géométriques
correspondants, et $\widetilde{U}=\fprod{U}{\widetilde{X}}{X}$. On a,
avec les notations évidentes,
$$
(f^{\ast}F)_{\bar{x}}=H^{0}(\widetilde{X},\widetilde{f}^{\ast}\widetilde{F})\simeq
H^{0}(\widetilde{S},\widetilde{F}){\quoi}.
$$
Par suite il suffit de démontrer que
$\widetilde{g}:\widetilde{U}\rightarrow \widetilde{S}$ est
$0$-acyclique, ce qui impliquera
$$
(j_{\ast}g^{\ast}F)_{\bar{x}}\simeq
H^{0}(\widetilde{U},\widetilde{g}^{\ast}\widetilde{F})\simeq
H^{0}(\widetilde{S},\widetilde{F}){\quoi}.
$$
Puisque $\widetilde{g}$ est limite de morphismes lisses, il est
localement $0$-acyclique (\hyperrefext[XV][XV.thm2.1]{2.1}, \hyperrefext[XV][XV.def1.11]{1.11}), 
et par \hyperrefext[XV][XV.cor1.12]{1.12} il suffit de
démontrer que les fibres géométriques de $\widetilde{U}/\widetilde{S}$
sont $0$-acycliques, c'est-à-dire connexes et non-vides. Soit
$\bar{s}'$ un point géométrique de $\widetilde{S}$. Alors la fibre
$\widetilde{X}_{\bar{s}'}$ est régulière, connexe et non vide, et il
résulte de l'hypothèse que $\widetilde{U}_{\bar{s}'}$ est un ouvert
dense de $\widetilde{X}_{\bar{s}'}$, donc connexe et non-vide, d'où le
résultat.

\begin{enonce*}{Corollaire 3.2.1}\etiquette[3.2.1]{XVI.cor3.2.1}
Sous les conditions de \hyperref[XVI.prop3.2]{3.2}, le foncteur $X'\mapsto X'|U$ de la
catégorie des revêtements étales de $X$ dans la catégorie des
revêtements étales de $U$ est pleinement fidèle. En particulier, pour
tout groupe fini $G$, l'application de restriction
$$
H^{1}(X,G)\longrightarrow H^{1}(U,G)
$$
est injective. 
\end{enonce*}
  
Cela\pageoriginale résulte de \hyperref[XVI.prop3.2]{3.2}, via un argument standard que nous
allons, pour la commodité des références, expliciter en un 

\begin{enonce*}{Lemme 3.2.2}\etiquette[3.2.2]{XVI.lem3.2.2}
Soit $f:X\rightarrow Y$ un morphisme de schémas, et soit
$f^{\ast}:\Et(Y)\rightarrow \Et(X)$ le foncteur image inverse
$Y'\mapsto \fprod{Y'}{X}{Y}$ de la catégorie $\Et(Y)$ des revêtements
étales de $Y$ dans $\Et(X)$. Conditions équivalentes:
\begin{itemize}
\item[(i)] $f^{\ast}$ est fidèle (\resp pleinement fidèle).

\item[(ii)] Pour tout revêtement étale $Y'$ de $Y$, désignant par
  $f':X'\rightarrow Y'$ le morphisme déduit de $f$ par changement de
  base $Y'\rightarrow Y$, l'application $U\mapsto {f'}^{-1}(U)$ de
  l'ensemble des parties à la fois ouvertes et fermées des $Y'$ dans
  l'ensemble des parties à la fois ouvertes et fermées de $X'$ est
  injective (\resp bijective).

\item[(ii bis)] Avec les notations de {\rm(ii)}, pour tout faisceau
  d'ensembles constant $C$ sur $Y'$, l'application canonique
$$
H^{0}(Y',C)\longrightarrow H^{0}(X',f^{\ast}C)
$$
est injective (\resp bijective).

\item[(iii)] Avec les notations de {\rm(ii)}, l'application canonique
$$
\Gamma(Y'/Y)\longrightarrow \Gamma(X'/X)
$$
est injective (\resp bijective).
\end{itemize}
\end{enonce*}

La démonstration est laissée au lecteur (\cf \SGA 2, IX 3.1 et 3.2).

\begin{enonce*}[remark]{Remarques 3.2.3}\etiquette[3.2.3]{XVI.rems3.2.3}
On\pageoriginale peut considérablement généraliser \hyperref[XVI.prop3.2]{3.2}, tant en termes
de faisceaux d'ensembles, en affaiblissant les hypothèses sur $f$,
qu'en termes de faisceaux de groupes avec des conclusions sur les
$R^{q}g_{\ast}(f^{\ast}F)$ (\cf \SGA 2 XIV 1 et \SGA 1 XIII). Remarque
analogue pour \hyperref[XVI.cor3.4]{3.4} à \hyperref[XVI.cor3.6]{3.6}.
\end{enonce*}

« Rappelons »  aussi le théorème de pureté de Zariski-Nagata
sous sa forme relative:

\begin{enonce*}{Théorème 3.3}\etiquette[3.3]{XVI.thm3.3}
(Théorème de pureté relatif). Soit
$$
\xymatrix@=1.25cm{
U \ar[rr]^{j} \ar[rd]_{g}&     &X\ar[ld]^{f}\\
  & S   &
}
$$
un diagramme commutatif où $f$ est lisse, et où $j$ est une immersion
ouverte telle que pour tout $s\in S$, $X_s-U_s$ soit partout de
codimension $\geq 2$ dans $X_s$. Alors le foncteur $X'\mapsto X'|U$ de
la catégorie des revêtements étales de $X$ dans la catégorie des
revêtements étales de $U$ est une \emph{équivalence de catégories.}
\end{enonce*}

\noindent
\emph{Démonstration}. Que le foncteur soit pleinement fidèle a été vu
(\hyperref[XVI.cor3.2.1]{3.2.1}). Le fait qu'il soit \sisi{ess.}{essentiellement} surjectif va se déduire du théorème
classique. On se donne un revêtement étale $U'$ de $U$, il faut
montrer qu'il se prolonge en un revêtement étale $X'$ de
$X$. Réduction immédiate facile au cas où $U$ est rétrocompact dans
$X$, puis au cas où $S$ est affine et de type fini sur $\Spec
\ZZ$, $X$ un schéma connexe, et $f$ de type fini.

On\pageoriginale procède par récurrence sur $n=\dim S$, l'assertion
étant conséquence du théorème classique (\SGA 2 X 3.4) dans le cas
$n=0$. Ayant l'unicité, on peut appliquer la technique de descente de
\sisi{\SGA IX 4.7}{\SGA 1 IX 4.7} au normalisé $S\leftarrow \bar{S}$ pour se réduire au cas
$S$ (et donc $X$) normal ($\bar{S}\rightarrow S$ étant fini (\EGA IV
7.8.3 (iii) (vi)). De plus, quitte à agrandir $U$, on peut supposer
$U$ maximal parmi les ouverts au-dessus desquels $U'$ peut se
prolonger, et il faut prouver alors $U=X$. Sinon, soient $x$ un point
maximal de $X-U$, et $s=f(x)$. Tout revient à prouver qu'on peut
prolonger $U'$ au-dessus d'un voisinage de $x$. Nous pouvons supposer
que $\dim \Oo_{S,s}=n\geq 1$ et que le théorème est déjà démontré pour
la dimension $<n$. Posons $B=\Oo_{X,x}$ et soit $B'$ le normalisé de $B$
dans l'anneau des fonctions rationnelles sur le revêtement $U'$ de
$U$; donc $B'$ est une $B$-algèbre finie (\loccit). Il suffit
évidemment de démontrer que $B'$ est une $B$-algèbre étale.

Soient $A=\Oo_{S,s}$, $0\neq t\in \max A$, $A_{0}=A/tA$,
$B_{0}=B/tB$. Par hypothèse de récurrence, la restriction du
revêtement à $U_{0}=U\otimes_A A_{0}$ s'étend à un revêtement
étale de $\Spec B_{0}$, appelons-le $\Spec
B^{\ast}_{0}$. Quitte à remplacer $X$ par un revêtement étale
convenable d'un voisinage de $x$ (ce qui est loisible) on peut
supposer $\Spec B^{\ast}_{0}$ complètement décomposé sur $\Spec
B_{0}$. On est ainsi ramené au cas où la $B_{0}$-algèbre
$B'_{0}$ est complètement décomposée en dehors du point fermé de
$\Spec B_{0}$. Alors les hypothèses de \hyperrefext[XV][XV.lem3.3]{3.3} sont satisfaites,
d'où le résultat.

\begin{enonce*}{Corollaire 3.4}\etiquette[3.4]{XVI.cor3.4}
Mêmes hypothèses que dans \hyperref[XVI.thm3.3]{3.3}. Alors on a
$(R^{1}j_{\ast})g^{\ast}F=0$ pour chaque faisceau $F$ de groupes
ind-finis sur $S$.
\end{enonce*}

Réduction\pageoriginale comme d'habitude au cas $S$ noethérien, $f$ de
type fini, et $F$ constructible. Soit $0\rightarrow F\rightarrow G$
une injection de faisceaux de groupes ind-finis sur $S$, et $C=G/F$. Il
résulte de \hyperref[XVI.prop3.2]{3.2} que le foncteur $j_{\ast}g^{\ast}$ est
exact, donc que 
$j_{\ast}g^{\ast}G\rightarrow j_{\ast}g^{\ast}C$ est surjectif, donc
(\hyperrefext[XII][XII.prop3.1]{3.1}) que $(R^{1}j_{\ast})g^{\ast}F\rightarrow
(R^{1}j_{\ast})g^{\ast}G$ est injectif. On peut donc remplacer $F$ par
$G$, et on se ramène ainsi par \hyperrefext[IX][IX.prop2.14]{2.14} et \hyperrefext[VIII][VIII.prop5.5]{5.5}
au cas où $F=G$ est constant, à valeur un groupe fini ordinaire $G$.
%% Cette réduction n'est pas parfaitement claire dans le cas non abélien :
%% donner référence pour R^1 j_* p_*G = p_* R^1 j_*G
Soit $\bar{x}$ un
point géométrique de $X$, $\widetilde{X}$ le localisé strict de $X$ en
$\bar{x}$, et $\widetilde{U}=\fprod{\widetilde{X}}{U}{X}$. Alors il
résulte de \hyperref[XVI.thm3.3]{3.3} que l'on a
$$
(R^{q}j_{\ast}g^{\ast}F)_{\bar{x}}=H^{q}(\widetilde{U},G_{\widetilde{U}})\simeq
H^{q}(\widetilde{X},G_{\widetilde{X}}){\quoi},
$$
et ce dernier est nul, d'où le résultat.

\begin{enonce*}{Lemme 3.5}\etiquette[3.5]{XVI.lem3.5}
(Lemme d'Abhyankar relatif). Soit $(Y,X)$ un $S$-couple lisse de
  codimension $1$, tel que $Y$ soit défini par une équation $t=0$ dans
  $X$. Soient $U=X-Y$ et $V/U$ un revêtement principal galoisien
  d'ordre $n$ premier aux caractéristiques résiduelles de $Y$. Posons
  $X'=\Spec \Oo_X[z]/(z^n-t)$, $U'=\fprod{U}{X'}{X}$,
  $V'=\fprod{V}{X'}{X}$. Alors le revêtement $V'$ s'étend uniquement à
  un revêtement étale de $X'$.
\end{enonce*}

\noindent
\emph{Démonstration}. Réduction immédiate au cas $S$, $X$, $Y$ affines
et de type fini sur $\Spec \ZZ$. L'unicité est un cas
particulier de \hyperref[XVI.cor3.2.1]{3.2.1}. Nous allons déduire l'existence de \hyperrefext[XV][XV.cor1.12]{1.12}:
soient $\bar{s}$ un point géométrique de $S$ et $\bar{x}$ un point
géométrique de $Y$ au-dessus de $\bar{s}$. Soient
$\widetilde{f}:\widetilde{X}\rightarrow \widetilde{S}$ le morphisme
des localisés stricts correspondants\pageoriginale induit par $f$,
$\widetilde{U}=\fprod{U}{\widetilde{X}}{X}$, etc\ldots Par descente,
compte tenu de l'unicité, il suffit évidemment de démontrer que
$\widetilde{V}'$ s'étend en un revêtement étale de $\widetilde{X}'$,
c'est-à-dire (puisque $\widetilde{X}'$ n'admet pas de revêtement
non-trivial) que $\widetilde{V}'$ est un revêtement trivial de
$\widetilde{U}'$.

Soit
$\bar{U}=\varprojlim_m\widetilde{U}_{m}=\Spec \Oo_{\widetilde{U}}[t^{1/m}]$,
où $m$ parcourt l'ensemble des entiers $m$ premiers à la
caractéristique résiduelle de $\widetilde{X}$. Alors
$\bar{U}/\widetilde{U}$ est un « revêtement galoisien infini »  de
groupe abélien $G=\varprojlim_m \ZZ/m\ZZ$, et on voit
immédiatement qu'il suffit de démontrer que $\bar{U}$ est
$1$-asphérique pour l'ensemble $\LL\subset \PP$
complémentaire à la caractéristique résiduelle de $\widetilde{X}$.

Considérons le morphisme $\bar{f}:\bar{U}\rightarrow
\widetilde{S}$. Il est limite de morphismes lisses, donc localement
$1$-asphérique pour $\LL$ (\hyperrefext[XV][XV.thm2.1]{2.1}, \hyperrefext[XV][XV.def1.11]{1.11}
(ii)). D'après \hyperrefext[XV][XV.cor1.12]{1.12}, il suffit de démontrer que les fibres
géométriques de $\bar{U}/\widetilde{S}$ sont $1$-asphérique pour
$\LL$. Soit $\bar{s}$ un point géométrique de $\widetilde{S}$. La
fibre géométrique $\bar{U}_{\bar{s}}$ est limite des fibres
$(\widetilde{U}_m)_{\bar{s}}$. Par suite il suffit de démontrer que
chaque revêtement principal galoisien d'un
$(\widetilde{U}_m)_{\bar{s}}$, d'ordre premier aux \sisi{car. rés.}{caractéristiques résiduelles},
induit un revêtement trivial de $\bar{U}_{\bar{s}}$. En remplaçant
$\widetilde{X}$ par
$\widetilde{X}_m=\Spec \Oo_{\widetilde{X}}[t^{1/m}]$ (qui est
encore lisse au-dessus de $\widetilde{S}$), on se réduit au cas $m=1$,
c'est-à-dire, $\bar{U}_m=\widetilde{U}$. Soit $V$ un tel revêtement de
$\widetilde{U}_{\bar{s}}$. Or $\widetilde{X}_{\bar{s}}=Z$ est un
schéma strictement local régulier, et $W=\widetilde{U}_{\bar{s}}$ est
l'ouvert complémentaire à un sous-ensemble régulier de codimension
$1$. On peut donc appliquer le lemme d'Abhyankar sous sa forme usuelle
(\SGA 1 X 3.6) pour conclure que $V$ induit un
revêtement\pageoriginale $V_m$ sur $W_m$ qui s'étend en un revêtement
étale de $Z_m$, pour $m$ convenable. Mais $Z$ donc $Z_m$ étant
strictement local, il s'ensuit que ce revêtement est trivial, donc
$V_m$ est un revêtement trivial de $W_m$, \cqfd

\begin{enonce*}[remark]{Remarque 3.5.1}\etiquette[3.5.1]{XVI.rem3.5.1}
Lorsqu'on ne suppose pas l'ordre $n$ de $G$ premier aux
caractéristiques résiduelles de $X$, on peut encore prouver ceci: soit
$m$ le plus grand entier premier aux caractéristiques résiduelles de
$Y$ qui divise l'ordre de $G$, et supposons que pour chaque point
$s\in S$, le revêtement $V_s$ de $U_s=X_s-Y_s$ ait, en les points
maximaux de $Y_s$, des groupes d'inertie d'ordres premiers aux
caractéristiques résiduelles de $Y$. Alors la conclusion de \hyperref[XVI.lem3.5]{3.5} reste
valable en prenant $X'=\Spec \Oo_X[z]/(z^m-t)$. On peut aussi
prendre pour $m$, plus généralement, un multiple commun quelconque des
ordres des groupes d'inertie qu'on vient d'envisager, supposés
premiers aux caractéristiques résiduelles correspondantes.
\end{enonce*}

\begin{enonce*}{Corollaire 3.6}\etiquette[3.6]{XVI.cor3.6}
(Pureté cohomologique en dimension $1$). Soit $(Y,X)$ un $S$-couple
  lisse de codimension $1$, et soit $F$ un faisceau de groupes
  constructible et localement constant sur $X$, d'ordres premiers aux
  caractéristiques résiduelles de $X$. Alors avec les notations de
  $3.1$, $(R^{1}j_{\ast})j^{\ast}F$ est localement isomorphe comme
  faisceau d'ensembles pointés à $i_{\ast}(F_Y/\cint(F_Y))$, où
  $F_Y=F|Y$ et où $F_Y/\cint(F_Y)$ désigne le quotient de $F_Y$ par
  les opérations de $F_Y$ sur lui-même par automorphismes intérieurs.
\end{enonce*}

\begin{enonce*}[remark]{Remarque}
L'isomorphisme du corollaire \hyperref[XVI.cor3.6]{3.6} n'est pas canonique.
\end{enonce*}

\noindent
\emph{Démonstration}.\pageoriginale On se ramène immédiatement au cas
$X$ noethérien et connexe. Puisque l'assertion est locale sur $X$ pour
la topologie étale, nous pouvons supposer que $F$ est un faisceau
constant, à valeur un groupe fini ordinaire $G$ d'ordre $m$ premier
aux caractéristiques résiduelles de $X$, et de plus que
$𝛍_{mX}\simeq (\ZZ/m\ZZ)_X$. Soit
$U=X-Y$ et $U'=\Spec \Oo_U[t^{1/m}]$, qui est un revêtement principal
galoisien « essentiel »  de $U$ de groupe $\ZZ/m\ZZ$
puisque $𝛍_{mX}\simeq (\ZZ/m\ZZ)_X$. Les revêtements principaux de groupe $G$
de $U$ qui sont trivialisés sur $U'$ sont classifiés localement sur
$Y$ par $H=\Hom(\ZZ/m \ZZ,G)$ mod $\cint(G)$, qui est
un ensemble pointé isomorphe à l'ensemble sous-jacent de
$G/\cint(G)$. On obtient ainsi un morphisme du faisceau constant $H_Y$
sur $Y$ dans $(R^{1}j_{\ast})G_{U}$, et il suffit de démontrer qu'il
est bijectif, ce qui résulte immédiatement de \hyperref[XVI.lem3.5]{3.5}.

\begin{enonce*}{Théorème 3.7}\etiquette[3.7]{XVI.thm3.7}
\emph{(Pureté cohomologique)}. Soit $(Y,X)$ un $S$-couple lisse de
codimension $c>0$. Soit $F$ un faisceau abélien sur $X$ localement
isomorphe (pour la topologie étale) à un faisceau de la forme
$f^{\ast}(G)$, où $G$ est un faisceau de torsion sur $S$, premier aux
caractéristiques résiduelles (par exemple $F$ un faisceau localement
constant de groupes finis d'ordres premiers aux caractéristiques
résiduelles de $X$). Alors, avec les notations de $3.1$, on a
\begin{align*}
&\HH^{q}_Y(X,F)=(R^{q}i^{!})F=0\qquad\mbox{si}\quad q\neq 2c,\\
&\begin{cases}
 F\xrightarrow{\sim} j_{\ast}j^{\ast}F\quad\et &\tag{\ie}\\
(R^{q}j_{\ast})j^{\ast}F=0 &\hspace{1.3cm} \mbox{si}\quad q\neq 0,\quad 2c-1,
\end{cases}
\end{align*}
et\pageoriginale
$$
\HH^{2c}_Y(X,F)=i^{\ast}[(R^{2c-1}j_{\ast})j^{\ast}F]
$$
est un faisceau localement isomorphe à $i^{\ast}(F)$.
\end{enonce*}

\noindent
\emph{Démonstration}. L'assertion est locale sur $X$ pour la topologie
étale; on peut donc \footnote[1]{Du moins si $F$ était supposé
  localement constant. Le lecteur se convaincra que la démonstration
  qui suit s'applique aussi, essentiellement au cas général.} % syntaxe étrange
supposer $F$ constant, donc si on veut $F\simeq f^{\ast}(\ZZ/n\ZZ)_S$
pour le faisceau constant $(\ZZ/n\ZZ)_S$ sur $S$.

Traitons d'abord le cas $c=1$. (L'assertion pour
$(R^{1}j_{\ast})j^{\ast}F$, dans le cas $F$ localement constant tout
au moins, résulterait immédiatement de \hyperref[XVI.cor3.6]{3.6}, mais nous allons le
déduire de nouveau.) Rappelons que $(Y,X)$ est localement isomorphe au
couple $(\EE^{n-1}_S, \EE^n_S)$, et si l'on remplace $S$
par $\EE^{n-1}_S$, on se ramène au cas où
$X=\EE^{1}_S=\Spec \Oo_S[t]$ et où $Y$ est la section $t=0$ de
$X/S$. La situation est localement isomorphe à l'inclusion de la
section à l'infini dans l'espace projectif $\PP^{1}_S$. On peut
donc prendre $X=\PP^{1}_S$, $Y$ la section à l'infini, et
$U=\EE^{1}_S=X-Y$.

Examinons la suite spectrale de Leray
$$
(R^{p}f_{\ast})(R^{q}j_{\ast})(\ZZ/n)\Longrightarrow
(R^{p+q}g_{\ast})(\ZZ/n){\quoi}.
$$
Or l'aboutissement est nul pour $p+q>0$, d'après \hyperrefext[XV][XV.cor2.2]{2.2}. De plus, les
$R^{q}j_{\ast}$ sont concentrés sur $Y$ si $q>0$, qui s'envoie
isomorphiquement sur $S$ par $f$, et d'autre part
$j_{\ast}(F|U)\xleftarrow{\sim}F$ en vertu de \hyperref[XVI.prop3.2]{3.2}. Il s'ensuit que
$(R^{p}f_{\ast})(R^{q}j_{\ast})=0$ si $p$ et $q>0$, et que
$f_{\ast}(R^{q}j_{\ast})$ détermine le faisceau $(R^{q}j_{\ast})$ si
$q>0$, enfin $\sisi{\RR}{R}^{p}f_{\ast}j_{\ast}(F|U)=\sisi{\RR}{R}^{p}f_{\ast}(F)$.
La suite spectrale se réduit donc à des isomorphismes\pageoriginale
$$
f_{\ast}(R^{q}j_{\ast})(F|U)\xrightarrow{\sim}
(R^{q+1}f_{\ast})(F)\quad\mbox{\rm si}\quad q>0{\quoi}.
$$
Mais $(R^{q}f_{\ast})(F)$ se calcule fibre par fibre (\hyperrefext[XII][XII.cor5.2]{5.2}), donc
est nul si $q\neq 0$, $2$, et est un faisceau localement isomorphe à
$G$ si $q=2$, comme il résulte aisément de \hyperrefext[XII][XII.cor5.2]{5.2} et \hyperrefext[IX][IX.cor4.7]{4.7}. Donc
$(R^{q}j_{\ast})(F|U)=0$ si $q\neq 1$, et est un faisceau localement
isomorphe à $i^{\ast}(F)$ si $q=1$.

Le cas $c>1$ se traite maintenant facilement par récurrence: soit
$(Y,X)$ un $S$-couple lisse de codimension $c>1$. Il est clair que,
localement sur $X$, on peut trouver un sous-schéma $Z$ de $X$ de
codimension 1 qui contient $Y$ et qui est lisse au-dessus de $S$, de
sorte qu'on ait un diagramme commutatif
$$
\xymatrix@=1.25cm{
Y       \ar[rr]^{u} \ar[rd]_{i}&    & Z\ar[ld]^v            \\
                  & X &,
}
$$
où chaque couple est un $S$-couple lisse (on devrait appeler $(Y,Z,X)$ un
$S$-triple lisse), où la codimension de $(Z,X)$ est $1$, et celle de
$(Y,Z)$ est $c-1$.

Or on a la suite spectrale de foncteurs composés
$$
(R^{p}u^{!})(R^{q}v^{!})F\Longrightarrow (R^{p+q}i^{!})F{\quoi}.
$$
Par hypothèse de récurrence, $(R^{q}v^{!})F=0$ si $q\neq 2$ et est
localement isomorphe à $v^{\ast}(F)$ si $q=2$, d'où encore par
hypothèse de récurrence, appliquée à $(R^{q}v^{!})F$, que
$(R^{p}u^{!})(R^{q}v^{!})F$ est nul si $(p,q)\neq
(2c-2,2)$\pageoriginale et est localement isomorphe à
$u^{\ast}w^{\ast}(F)$ si $(p,q)=(2c-2,2)$, d'où immédiatement le
résultat pour $(R^{p+q}i^{!})F$, \cqfd

\begin{enonce*}{Corollaire 3.8}\etiquette[3.8]{XVI.cor3.8}
Soient $(Y,X)$ un $S$-couple lisse de codimension $c$, $n$ un entier
naturel premier aux caractéristiques résiduelles de $X$, et posons
$$
T_{Y/X}=\HH^{2c}_Y(\ZZ/n\ZZ){\quoi},
$$
qui est un faisceau sur $Y$ localement isomorphe (pour la topologie
étale) au faisceau $(\ZZ/n\ZZ)_Y$ en vertu de
$(3.7)$. Soit $F$ un faisceau abélien sur $X$, annulé par $n$,
satisfaisant à la condition énoncée dans $3.7$. Alors on a (pour
mémoire) $\HH^{i}_Y(F)=0$ pour $i\neq 2c$, et de plus on a un
isomorphisme canonique:
$$
\HH^{2c}_Y(F)\simeq i^{\ast}(F)\otimes T_{Y/X}{\quoi}.
$$
\end{enonce*}

\noindent
\emph{Démonstration}. On définit aisément un homomorphisme canonique
$$
\varphi:i^{\ast}(F)\otimes T_{Y/X}\longrightarrow \HH^{2c}_Y(F)
$$
\ie un homomorphisme
$$
i^{\ast}(F)\longrightarrow
\sheaf{\Hom}(\HH^{2c}_Y(\ZZ/n\ZZ),H^{2c}_Y(F)){\quoi},
$$
en utilisant l'homomorphisme canonique $F\rightarrow
\sheaf{\Hom}(\ZZ/n\ZZ,F)$. Il reste à prouver que $\phi$ est
un isomorphisme, ce qu'on vérifie aisément en suivant la démonstration
donnée de \hyperref[XVI.thm3.7]{3.7}.

\begin{enonce*}{Corollaire 3.8.1}\etiquette[3.8.1]{XVI.cor3.8.1}
Sous les conditions de \hyperref[XVI.cor3.8]{3.8}, la formation des
$\HH^{i}_Y(F)$\pageoriginale commute à tout changement de base
$S'\rightarrow S$.
\end{enonce*}

C'est clair.

\begin{enonce*}{Corollaire 3.9}\etiquette[3.9]{XVI.cor3.9}
(Théorème de pureté cohomologique absolue). Soient $X$ un schéma
régulier, $Y$ un sous-préschéma fermé régulier de codimension $c$ en
chaque point. Supposons de plus que $X$ soit localement de type fini
sur un corps parfait $k$. Alors, si $n$ est premier aux
caractéristiques résiduelles de $X$, on a
$$
\HH^{i}_Y(\ZZ/n\ZZ)=0\mbox{ si } i\neq 2c{\quoi},
$$
et $H^{2c}_Y(\ZZ/n\ZZ)$ est un faisceau sur $Y$
localement isomorphe (pour la topologie étale) au faisceau constant
$\ZZ/n\ZZ$). 
\end{enonce*}

C'est un cas particulier de \hyperref[XVI.thm3.7]{3.7}, compte tenu que l'hypothèse implique
que $X$ et $Y$ sont lisses sur $k$.

\begin{enonce*}[remark]{Remarques 3.10}\etiquette[3.10]{XVI.rems3.10}
\begin{enumerate}
\renewcommand{\theenumi}{\alph{enumi}}
\renewcommand{\labelenumi}{\theenumi)}
\item On peut dans \hyperref[XVI.cor3.8]{3.8} expliciter la structure du faisceau $T_{Y/X}$,
on trouve qu'il est canoniquement isomorphe au faisceau
$(𝛍_n)_Y^{\otimes -c}$, où $𝛍_n$
désigne le faisceau des racines $n$.èmes de l'unité (localement
isomorphe au faisceau constant
$\ZZ/n\ZZ$). L'isomorphisme en question est étudié sous
le nom de \emph{classe fondamentale} locale de $Y$ dans $X$, dans la
situation un peu plus générale des intersections complètes relatives,
dans Exp. XVIII consacré à la dualité.

\item Il\pageoriginale est très plausible que les conclusions de
\hyperref[XVI.cor3.9]{3.9} et de \hyperref[XVI.rems3.10]{3.10} a) restent valables sans
postuler l'existence d'un corps de base $k$, tout au moins lorsque $X$
est un préschéma « excellent »\nde{Travaux de Gabber XXX}. C'est ce qui sera établi en tous
cas dans Exp. XIX lorsque $X$ est de caractéristique nulle, en
utilisant la résolution des singularités de Hironaka. Comme nous
verrons dans \SGA 5 on a besoin du « théorème de pureté absolu »
(ainsi que de la résolution des singularités) notamment pour pouvoir
établir la « formule de dualité locale », (qui elle-même est un
des ingrédients majeurs de la formule de \textsc{Lefschetz-Verdier}).
\end{enumerate}
\end{enonce*}

\section{Théorème de comparaison de la cohomologie pour les préschémas
algébriques sur $\CC$.}\etiquette[4]{XVI.sec4}

Soit $X$ un schéma algébrique sur $\Spec \CC$. Nous allons
comparer les sites étale et classique sur $X$, et nous reprenons les
notations de \hyperrefext[XI][XI.sec4]{4}. Le théorème de comparaison, qui est la forme
générale de \hyperrefext[XI][XI.thm4.4]{4.4}, est le suivant:

\begin{enonce*}{Théorème 4.1}\etiquette[4.1]{XVI.thm4.1}
Soit $f:X\rightarrow S$ un morphisme de type fini de schémas
localement de type fini sur $\Spec \CC$, de sorte qu'on a un
diagramme commutatif de morphismes de sites
$$
\xymatrix@C=2cm@R=1.5cm{
X_{\et}  \ar[d]_{f_{\et}}                   &   X_{\cl}
\ar[l]_{\varepsilon} \ar[d]^{f_{\cl}} \\
S_{\et}                    &   S_{\cl}\ar[l]_{\varepsilon} 
}
$$
Soit\pageoriginale $F$ un faisceau d'ensembles (\resp de groupes
ind-finis, \resp abélien de torsion) et supposons qu'on soit dans l'un
des cas suivants:
\begin{enumerate}
\renewcommand{\theenumi}{\roman{enumi}}
\renewcommand{\labelenumi}{(\theenumi)}
\item $f$ propre,

\item $F$ constructible.
\end{enumerate}
Alors les morphismes de changement de base \sisi{\emph{XII 4.1.2}}{\hyperrefext[XII][XII.eq4.2]{4.2}}
\footnote[1]{En fait, nous utilisons l'homomorphisme de changement de
base sous des conditions plus générales que celles envisagées dans
Exp. XII, où nous nous étions limités à un carré \emph{cartésien} de
morphismes de \emph{préschémas}. Le lecteur se convaincra aisément que
la définition s'étend verbatim à un carré \emph{essentiellement
commutatif} de morphismes de \emph{sites}. Pour plus de développements
à ce sujet, voir exposé suivant.}
$$
\varphi:\epsilon^{\ast}(R^{q}f_{\et\ast})F\longrightarrow
(R^{q}f_{\cl\ast})\epsilon {}^{\ast}F
$$
sont bijectifs pour $q=0$ (\resp pour $q=0$, $1$, \resp pour $q\geq
0$).
\end{enonce*}

Remarquons que le théorème de Grauert-Remmert (\hyperrefext[XI][XI.thm4.3]{4.3} (iii)) que nous
avions utilisé dans la démonstration de \hyperrefext[XI][XI.thm4.4]{4.4}, est un cas particulier
de \hyperref[XVI.thm4.1]{4.1}.

Pour\pageoriginale la démonstration, nous traitons d'abord le cas $f$
propre. Dans ce cas, le résultat sera conséquence de GAGA et du
théorème de changement de base pour un morphisme propre, par un
raisonnement élémentaire qui aurait pu figurer dans Exp. XIV. Pour
traiter le cas général, nous utiliserons de plus la résolution des
singularités [1] et \hyperref[XVI.thm3.7]{3.7}.

On peut évidemment supposer $Y$, donc aussi $X$, de type fini sur
$\CC$.

\emph{Cas $f$ propre.} Nous rappelons les faits élémentaires suivants:
\begin{enumerate}
\renewcommand{\theenumi}{\alph{enumi}}
\renewcommand{\labelenumi}{(\theenumi)}
\item L'espace topologique $X_{\cl}$ est localement compact.

\item Si $f:X\rightarrow S$ est propre, alors
$f_{\cl}:X_{\cl}\rightarrow S_{\cl}$ est une application \emph{propre}
d'espaces topologiques.
\end{enumerate}

En effet, (a) est trivial puisque localement $X_{\cl}$ est un fermé
dans un $\CC^n$. Pour (b) on se ramène par le lemme de Chow au
cas où $f$ est projectif, ce cas étant trivial.

Soit $F$ un faisceau d'ensembles (\resp\ldots) sur $X_{\et}$. 
Pour les deux morphismes $f_{\et}$ et $f_{\cl}$, la
formation des $R^{q}$ commute à la formation des fibres (\hyperrefext[XII][XII.cor5.2]{5.2} et
\cite{SGA4-16:2} 4.11.1). Il suffit donc de vérifier le théorème fibre par fibre,
en les points fermés de $S_{\et}$.

Dans le lemme suivant, nous ne supposons pas que $g$ soit propre:

\begin{enonce*}{Lemme 4.2}\etiquette[4.2]{XVI.lem4.2}
Soit\pageoriginale
$$
\xymatrix@=1.5cm{
X  \ar[rd]_{g}& &\bar{X} \ar[ll]_{h}\ar[ld]^{\bar{g}}\\
& S &
}
$$
un diagramme commutatif avec $h$ propre, et soit $U$ un ouvert de $X$
tel que $h$ induise un isomorphisme de l'ouvert
$\bar{U}=\fprod{U}{\bar{X}}{X}$ de $\bar{X}$ sur $U$. Soient
$j:U\rightarrow X$ et $\bar{j}:\bar{U}\rightarrow \bar{X}$ les
inclusions, $F$ un faisceau d'ensembles (\resp\ldots) sur $U$, et
$\bar{F}$ le faisceau induit sur $\bar{U}$. Si le théorème $4.1$ est
vrai pour $(\bar{g},\bar{j}_{!}\bar{F})$ il l'est également pour $(g,j_{!}F)$.
\end{enonce*}

\noindent
\emph{Démonstration}. Le théorème est vrai pour
$(\bar{h},\bar{j}_{!}\bar{F})$ parce qu'on est ramené à le vérifier
fibre par fibre, et on a alors les deux cas suivants: ou bien la
restriction du faisceau à la fibre est nulle, ou bien le morphisme $h$
est un isomorphisme. Dans ces deux cas, le résultat est trivial. On
trouve que
$$
\epsilon^{\ast}j_{!}F\simeq \epsilon {}^{\ast}h_{\text{\rm
et\,}\ast}(\bar{j}_{!}\bar{F})\simeq
h_{\cl\ast}\epsilon^{\ast}(\bar{j}_{!}\bar{F})
$$
et
$$
(R^{q}h_{\text{\rm
et\,}\ast})\bar{j}_{!}\bar{F}=(R^{q}f_{\cl\ast})\epsilon^{\ast}(\bar{j}_{!}\bar{F})=0\text{ si } q>0
$$
pour les valeurs de $q$ envisagées, donc que
$$
(R^{q}g_{\text{\rm et\,}\ast})j_{!}F\simeq (R^{q}\bar{g}_{\text{\rm
et\,}\ast})\bar{j}_{!}\bar{F}
$$
et 
$$
(R^{q}g_{\cl \ast})j_{!}F\simeq (R^{q}\bar{g}_{\cl
\ast})\bar{j}_{!}\bar{F}{\quoi},
$$
d'où 4.2.\pageoriginale

\begin{enonce*}{Lemme 4.3}\etiquette[4.3]{XVI.lem4.3}
Pour vérifier \hyperref[XVI.thm4.1]{4.1} pour toutes les données $(f,F)$ avec $f$ propre,
il suffit de le faire dans le cas où on suppose que de plus $f$ est de
dimension relative $\leq 1$ et $S=\Spec \CC$ est un point.
\end{enonce*}

\noindent
\emph{Démonstration}. Récurrence sur la dimension relative. Supposons
le résultat connu en dimension relative $<n$, et que $n>1$. Comme nous
l'avons remarqué, nous pouvons faire la vérification fibre par fibre,
c'est-à-dire, nous pouvons supposer que $S=\Spec \CC$ est un
point, et donc $X$ de dimension $\leq n$. De plus, nous pouvons
supposer $X$ réduit. Soient $\varphi$ une fonction rationnelle sur $X$
qui n'est constante sur \sisi{aucun composant}{aucune composante} irréductible de $X$, $U$ un
ouvert dense de $X$ sur lequel $\varphi$ est défini, et
$\bar{X}\subset X\times \PP^{1}$ l'adhérence du graphe du
morphisme $\varphi:U\rightarrow \EE^{1}$. On a un diagramme
commutatif
$$
\xymatrix@R=1.5cm@C=1.5cm{
&\bar{X}\ar[ld]_{h} \ar[rd]^{a} \ar[dd]^{\bar{g}}   &\\
X  \ar[rd]_{g}  &                                  &\PP^1 \ar[ld]_{b}\\
            & **[r]Y              =\Spec \CC&
}
$$
où les dimensions relatives de $a$ et $b$ sont $<n$. Par l'hypothèse
de récurrence, le théorème est vrai pour $(a,.)$ et $(b,.)$, d'où on
conclut qu'il l'est également pour $(\bar{g},.)$. En effet, c'est
trivial pour $q=0$; si $F$ est abélien de torsion sur
$\bar{X}$, cela résulte\pageoriginale du morphisme de suites
spectrales de Leray
$$
\xymatrix{
\varepsilon^{\ast}(R^{p}h_{\text{et}\ast})(R^{q}a_{\text{et}\ast})F \ar@{=>}[r] \ar[d]&
\varepsilon^{\ast}(R^{p+q}\bar{g}_{\text{et}\ast})F  \ar[d]\\
(R^{p}b_{\text{cl}\ast}(R^{q}{a_{\text{cl}\ast}}))
\varepsilon^{\ast}F  \ar@{=>}[r] & (R^{p+q}\bar{g}_{\text{cl}\ast})\varepsilon^{\ast}F;
}
$$
si enfin $F$ est un faisceau de groupes ind-finis sur $\bar{X}$,
notons que puisque $\epsilon^{\ast}(R^{1}\bar{g}_{\text{\rm
et\,}\ast})$ est un foncteur effaçable, il suffit de démontrer
(compte tenu du résultat pour $q=0$), que
$(R^{1}\bar{g}_{\cl\ast})\epsilon^{\ast}$ est également effaçable
(\cf \hyperrefext[XII][XII.cor8.2]{8.2}), ce qui résulte de la suite exacte
$$
0\longrightarrow (R^{1}b_{\cl\ast})a_{\cl
\ast}\epsilon^{\ast}F\longrightarrow (R^{1}\bar{g}_{\cl
\ast})\epsilon^{\ast}F\longrightarrow b_{\cl \ast}(R^{1}a_{\cl
\ast})\epsilon^{\ast}F{\quoi},
$$
et du fait que $(R^{1}b_{\cl\ast})\epsilon^{\ast}$ et
$(R^{1}a_{\cl\ast})\epsilon^{\ast}$ sont effaçables grâce à
l'hypothèse de récurrence.

Soit maintenant $F$ un faisceau d'ensembles sur $X$. On a
$F\hookrightarrow h_{\text{\rm et\,}\ast}h_{\et}{}^{\ast}F=G$. Par la suite exacte
\begin{equation*}
F\rightarrow G\rightrightarrows G\amalg_F G{\quoi},\tag{4.3.1}
\end{equation*}\etiquette[4.3.1]{XVI.eq4.3.1}
on se ramène (pour $q=0$) à démontrer \hyperref[XVI.thm4.1]{4.1} pour $G$, donc pour un
faisceau de la forme $h_{\text{\rm et\,}\ast}\bar{F}$. Mais pour un
tel faisceau, le résultat sera conséquence du théorème pour
$(\bar{g},.)$ et pour $(h,.)$, et $h$ est de dimension relative $<n$,
comme on voit immédiatement.

Soit $F$ un faisceau abélien de torsion sur $X$. On a la suite exacte
\begin{equation*}
0\longrightarrow j_{!}j^{\ast}F\longrightarrow F\longrightarrow
i_{\ast}i^{\ast}F\longrightarrow 0{\quoi}.\tag{$\ast$}
\end{equation*}
Puisque $U$ est dense dans $X$, on a $\dim Y\leq n$, donc \hyperref[XVI.thm4.1]{4.1} est 
vrai\pageoriginale pour $(g|Y, i^{\ast}F)$, donc pour
$(g,i_{\ast}i^{\ast}F)$. Il suffit donc, grâce au lemme des cinq, de
vérifier \hyperref[XVI.thm4.1]{4.1} pour $(f,j_{!}j^{\ast}F)$, ce qui résulte du théorème
pour $\bar{g}$ (\hyperref[XVI.lem4.2]{4.2}).

Soit $F$ un faisceau de groupes ind-finis sur $X$. Il suffit de
démontrer que le foncteur $(R^{1}g_{\cl\ast})\epsilon^{\ast}$ est
effaçable pour le faisceau $F$, ce qui résulte de ($\ast$) et de
l'effaçabilité pour $j_{!}j^{\ast}F$ et pour $i_{\ast}i^{\ast}F$,
\ie du théorème pour $(g,j_{!}j^{\ast}F)$ (\hyperref[XVI.lem4.2]{4.2}) et pour
$(g,i_{\ast}i^{\ast}F)$ (hypothèse de récurrence), \cqfd

\begin{enonce*}{Lemme 4.4}\etiquette[4.4]{XVI.lem4.4}
Pour démontrer $(4.1)$ dans le cas où $f$ est propre, il suffit de
démontrer ceci: soit $X$ une courbe complète régulière sur $\Spec
\CC$, et $F$ un faisceau constant et constructible d'ensembles
(\resp\ldots) sur $X$. Alors les morphismes
$$
H^{q}(X_{\et},F)\longrightarrow H^{q}(X_{\cl},F)
$$
sont bijectifs pour les valeurs de $q$ envisagées.
\end{enonce*}

\noindent
\emph{Démonstration}. Puisqu'il suffit de vérifier le théorème fibre
par fibre, on est ramené par \hyperref[XVI.lem4.3]{4.3} au cas $X$ propre et de dimension
$\leq 1$. Le cas de dimension $0$ est d'ailleurs trivial, et il
résulte de cela que le théorème est vrai pour un morphisme
fini. Puisque $f$ est propre, la cohomologie de $X_{\et}$ et
de $X_{\cl}$ commute aux limites inductives (\hyperrefext[VII][VII.prop3.3]{3.3}
et [2], 4.12.1) et on est donc ramené au cas $F$ constructible (\hyperrefext[IX][IX.prop2.9]{2.9} (iii)).

On applique \hyperrefext[IX][IX.prop2.14]{2.14}. Pour un faisceau d'ensembles, la suite exacte
\eqhyperref[XVI.eq4.3.1]{4.3.1} et \hyperrefext[IX][IX.prop2.14]{2.14} montrent qu'on peut prendre $F=\pi{}_{\ast}C$,
où\pageoriginale $\pi:X'\rightarrow X$ est fini, $X'$ est normal, donc
régulier, et $C$ est un faisceau constant sur $X'$. Pour un faisceau
abélien de torsion, on prend une résolution de $F$ par des produits
finis de faisceaux de la forme $\pi_{\ast}C$ et on applique la suite
spectrale d'une résolution, et on se réduit encore au cas
$F=\pi_{\ast}C$. Pour un faisceau de groupes ind-finis, on rappelle
qu'il suffit de démontrer l'effaçabilité de
$(R^{1}g_{\cl\ast})\epsilon^{\ast}$ pour le faisceau $F$, donc pour un
faisceau $G$ dans lequel on peut plonger $F$, donc on est encore
réduit au cas $F=\pi_{\ast}C$. Or il est évident par (\hyperrefext[VIII][VIII.prop5.5]{5.5} et \hyperrefext[VIII][VIII.rem5.8]{5.8})
qu'on peut maintenant remplacer $X$ par $X'$, d'où le lemme.

Pour traiter ce dernier cas, il suffit évidemment (tenant compte de la
dimension cohomologique (\sisi{IX 7.7}{\hyperrefext[IX][IX.cor5.7]{5.7}}
et \cite{SGA4-16:2} 4.14.1)) de démontrer que pour
une courbe complète et régulière on a
\begin{itemize}
\item[$(q=0)$:] $X$ connexe et non-vide $\Leftrightarrow X_{\cl}$
connexe et non-vide.

\item[$(q=1)$:] Le morphisme de changement de base induit une
équivalence de la catégorie des revêtements étales de $X$ et de la
catégorie des revêtements étales analytiques finis de $X_{\cl}$.

\item[$(q=2)$:] On a un isomorphisme
$\epsilon:H^{2}(X_{\et},𝛍_n)\xrightarrow{\sim}
H^{2}(X_{\cl},𝛍_n)$.
\end{itemize}

Or les deux premières assertions résultent immédiatement de GAGA. Pour
la dernière, notons qu'on a $H^{2}(X_{\et},\GG_m)=0$ (\hyperrefext[IX][IX.eq4.5]{4.5})
et de même, notant $\Oo^{\ast}$ le faisceau des fonctions
holomorphes inversibles de $X_{\cl}$, $H^{2}(X_{\cl},\Oo^{\ast})=0$
(cela résulte de la suite exacte
$$
0\longrightarrow \ZZ\longrightarrow \Oo\xrightarrow{\exp}
\Oo^{\ast}\longrightarrow 0
$$
et\pageoriginale du fait que $H^{2}(X_{\cl},\Oo)=H^{2}(X_{\et},\Oo_X)=0$ par GAGA). Par suite la suite exacte de Kummer \hyperrefext[IX][IX.Thr3.2]{3.2}
(\resp la suite exacte sur $X_{\cl}$
$$
0\rightarrow 𝛍_n\rightarrow
\Oo^{\ast}\xrightarrow{n}\Oo^{\ast}\rightarrow 0)
$$
donne 
\begin{align*}
& H^{1}(X_{\et}, \GG_m)/n\xrightarrow{\sim}
H^{2}(X_{\et},𝛍_n)\\
(\resp & H^{1}(X_{\cl},\Oo^{\ast})/n\xrightarrow{\sim}
H^{2}(X_{\cl},𝛍_n)){\quoi}.
\end{align*}
D'après GAGA, le morphisme canonique
$$
H^{1}(X_{\et},\GG_m)=\Pic X\longrightarrow \Pic
X_{\cl}=H^{1}(X_{\cl},\Oo^{\ast})
$$
est bijectif, d'où le résultat, ce qui prouve \hyperref[XVI.thm4.1]{4.1} dans le cas où $f$
est propre.

\medskip
\noindent
\emph{Cas $F$ constructible}. La démonstration se fait par récurrence
sur $\dim X$. Supposons le résultat connu si la dimension est $<n$, et
prouvons-le quand la dimension de $X$ est $\leq n$.

\begin{enonce*}{Lemme 4.5}\etiquette[4.5]{XVI.lem4.5}
On peut supposer $f:X\rightarrow S$ une immersion ouverte dense et $F$
constant.
\end{enonce*}

\noindent
\emph{Démonstration}. Comme dans la démonstration de \hyperref[XVI.lem1.3]{1.3}, on peut en
effet supposer $X$ et $S$ affines, donc $f$ quasi-projectif. Alors on
applique le raisonnement de \hyperref[XVI.lem4.4]{4.4} pour se ramener au cas
$F=\pi_{\ast}C$, $\pi:X'\rightarrow X$ fini, et $C$ constant, d'où en
remplaçant $X$ par $X'$, au\pageoriginale cas $F$ constant. 
On peut remplacer $X$ par $X_{\red}$. Soit
$$
\xymatrix@=1.25cm{
  X  \ar[rr]^{h}\ar[rd]_{f} &      &\bar{X} \ar[ld]^{\bar{f}}\\
             &     S  &   
}
$$
un diagramme commutatif tel que $\bar{X}$ soit réduit, que $\bar{f}$
soit projectif, et que $h$ soit  une immersion ouverte dense. Il
suffit de vérifier le théorème pour les deux morphismes $h$, $\bar{f}$
(\cf dém. de \hyperref[XVI.lem4.3]{4.3}), donc pour $h$ puisque $\bar{f}$ est propre, donc
pour une immersion ouverte dense.

\begin{enonce*}{Lemme 4.6}\etiquette[4.6]{XVI.lem4.6}
On peut de plus supposer $S$ régulier.
\end{enonce*}

\noindent
\emph{Démonstration}. Soit $S'\xrightarrow{h}S$ une « résolution
des singularités de $S$ », \ie, un morphisme surjectif, propre et
birationnel, tel que $S'$ soit régulier, et soit
$$
\xymatrix@C=2cm@R=1.25cm{
X'     \ar[r]^{f'} \ar[d]^{h'}&     S'\ar[d]^{h}\\
X   \ar[r]^{f}  & S
}
$$
le diagramme cartésien déduit de $h$.

Soit $F$ un faisceau d'ensembles constant sur $X$. On a
$F\hookrightarrow G=h'_{\ast}{h'}^{\ast}F$ et on en déduit qu'il
suffit (pour $q=0$) de vérifier \hyperref[XVI.thm4.1]{4.1} pour $G$ (\cf \eqhyperref[XVI.eq4.3.1]{4.3.1}),
c'est-à-dire, pour un faisceau de la forme ${h'}_{\ast}F'$, où $F'$
est constant sur $X'$. Or on a un diagramme\pageoriginale commutatif
(où on supprime le symbole « et »)
$$
\xymatrix@C=2cm@R=1.5cm{
\varepsilon^{\ast}f_{\ast} (h'_{\ast}F') = \varepsilon^{\ast}h_{\ast}f'_{\ast}F'  \ar[d]^{d}\ar[r]^(.6){a} & h_{\text{cl}\ast}
	    \varepsilon^{\ast}f_{\ast}F' \ar[d]^{b}\\
	    f_{\text{cl}\ast}\varepsilon^{\ast}(h'_{\ast}F')  \ar[r]^(.4){c} &
	    f_{\text{cl}\ast} h'_{\text{cl}\ast}\varepsilon^{\ast}F' = h_{\text{cl}\ast}f'_{\text{cl}\ast}
	    \varepsilon^{\ast}F',
}
$$
et $a$, $c$ sont bijectifs d'après le résultat pour un morphisme
propre. Pour démontrer que $d$ est bijectif, il suffit de démontrer
que $b$ l'est, donc que
$$
\epsilon^{\ast}f'_{\ast}F'\xrightarrow{\sim}
f'_{\cl\ast}\epsilon^{\ast}F'
$$
donc on est ramené au cas $S=S'$, donc $S$ régulier.

\noindent
(Remarque: \sisi{Nous}{nous} avons utilisé sans le mentionner explicitement le
théorème de finitude pour un morphisme propre \hyperrefext[XIV][XIV.thm1.1]{1.1}).

Soit $F$ un faisceau abélien constant de torsion sur $X$. Il suffit de
traiter le cas $F=(\ZZ/n)_X$. Or le morphisme $h'$ induit un
isomorphisme au-dessus d'un ouvert dense $j:U\rightarrow X$ de $X$;
soient $j':U'\rightarrow X'$ l'ouvert
${h'}^{-1}(U)\xrightarrow{\sim}U$, et $Y=X-U$. On a $\dim Y<n$ parce
que $U$ est dense, donc le théorème est vrai pour
$(f,(\ZZ/n)_Y)$. Par la suite exacte
$$
0\longrightarrow j_{!}(\ZZ/n)_U\longrightarrow
(\ZZ/n)_X\longrightarrow (\ZZ/n)_Y\longrightarrow 0
$$
et le lemme des cinq, on se ramène à démontrer le théorème pour le
faisceau $j_{!}(\ZZ/n)_U$ et le morphisme $f$, donc par \hyperref[XVI.lem4.2]{4.2} à
démontrer le théorème pour $(hf', {j'_!}(\ZZ/n)_{U'})$. Or on
a une suite exacte
$$
0\longrightarrow {j'_!}(\ZZ/n)_{U'}\longrightarrow
(\ZZ/n)_{X'}\longrightarrow (\ZZ/n)_{Y'}\longrightarrow
0{\quoi},
$$
et\pageoriginale le théorème est vrai pour $(hf',(\ZZ/n)_{Y'})$
par l'hypothèse de récurrence. Il suffit donc de démontrer le théorème
pour $(hf',(\ZZ/n)_{X'}$). Mais par le cas propre, le théorème
est vrai pour $(h,.)$, et il s'ensuit, par la suite spectrale de Leray
pour le couple de morphismes $h$, $f'$, qu'il suffit de le démontrer
pour $(f',\ZZ/n)$, d'où le lemme dans ce cas.

Supposons enfin que $F$ soit un faisceau de groupes finis
constant. Rappelons qu'il suffit de démontrer que le foncteur
$(R^{1}f_{\cl \ast})\epsilon^{\ast}(\cdot)$ est effaçable pour $F$,
donc pour un $G$ dans lequel on peut plonger $F$, donc pour
$G={h'_\ast}{h'}^{\ast}F$, c'est-à-dire pour un faisceau de la forme
${h'_\ast}F'$ où $F'$ est constant sur $X'$. On a
$$
\epsilon^{\ast}({h'_\ast}F')\xrightarrow{\sim}h'_{\cl \ast}(\epsilon^{\ast}F')
$$
et
$$
0\longrightarrow
(R^{1}f_{\cl\ast})h'_{\cl\ast}\epsilon^{\ast}(\cdot)\longrightarrow
(R^{1}f_{\cl}h'_{\cl\ast})\epsilon^{\ast}(\cdot)=(R^{1}h_{\cl}f'_{\cl\ast})\epsilon^{\ast}(\cdot){\quoi},
$$
donc il suffit de démontrer l'effaçabilité de ce dernier foncteur
pour un faisceau constant. Mais on a la suite exacte
$$
0\longrightarrow (R^{1}h_{\cl \ast})f'_{\cl
\ast}\epsilon^{\ast}(\cdot)\longrightarrow (R^{1}h_{\cl}f'_{\cl
\ast})\epsilon^{\ast}(\cdot)\longrightarrow
h_{\cl\ast}(R^{1}f'_{\cl\ast})\epsilon^{\ast}(\cdot){\quoi},
$$
et le premier membre est effaçable, disons effacé par la résolution
de Godement (\hyperrefext[XII][XII.prop3.3]{3.3}), parce que $h$ est propre. Il suffit donc de
démontrer que le membre de droite est effacé par la résolution de
Godement, donc que $(R^{1}f'_{\cl\ast})\epsilon^{\ast}(\cdot)$ l'est,
donc que le théorème est vrai pour $(f',F')$ où $F'$ est constant,
d'où le lemme.

\noindent
\emph{Fin de la démonstration}.\pageoriginale Nous supposons 
maintenant que $F$ est
un faisceau constant et que $f:X\rightarrow S$ est une immersion
ouverte dense, avec $S$ régulier. Le cas ensembliste résulte
maintenant immédiatement de \hyperref[XVI.prop3.2]{3.2} et du lemme trivial analogue pour la
topologie classique\sisi{ ---~en effet}{. En effet}, on en déduit que
$\epsilon^{\ast}f_{\ast}F$ et $f_{\cl\ast}\epsilon^{\ast}F$ sont
constants de même valeur, donc isomorphes par $\varphi$.

Pour $q>0$, on va d'abord se réduire au cas où de plus $X=S-Y$ avec
$Y$ régulier, \ie au cas où $(Y,S)$ est un $\Spec \CC$-couple lisse.
En effet, écrivons $X=S-C_1$, où $C_1$ est muni de la
structure induite réduite, et définissons récursivement

$Y_{\nu}=$ l'ouvert dense des points réguliers de $C_{\nu}$,

$C_{\nu +1}=C_{\nu}-Y_{\nu}$, avec structure induite réduite.

\noindent
On a $\dim C_{\nu+1}<\dim C_{\nu}$ parce que $Y_{\nu}$ est dense dans
$C_{\nu}$, d'où une suite
$$
X=X_1\subset X_2=S-C_2\subset X_3=S-C_3\subset\ldots \subset
X_r=S{\quoi},
$$
où les inclusions $i_{\nu}:X_{\nu}\rightarrow X_{\nu +1}$ sont des
immersions ouvertes et où l'on a
$$
X_{\nu}=X_{\nu+1}-Y_{\nu}, Y_{\nu}\;\mbox{\rm fermé dans
$X_{\nu+1}$, et régulier}{\quoi}.
$$ Or l'image directe $i_{\nu\ast}F$ d'un faisceau constant sur $X$
est constante (\hyperref[XVI.prop3.2]{3.2}), et les $R^{q}i_{\nu\ast}F$ sont
constructibles et concentrés sur les variétés de $Y$ de dimension
$\leq n-1$ (\hyperref[XVI.cor3.4]{3.4}, \hyperref[XVI.cor3.6]{3.6}, \hyperref[XVI.thm3.7]{3.7}) pour
les valeurs de $q$ envisagées. Donc par l'hypothèse de récurrence et
la suite spectrale de Leray (\resp la suite exacte \hyperrefext[XII][XII.prop3.2]{3.2}), on se
réduit\pageoriginale au cas de $i_{\nu}$, d'où la réduction annoncée.

Mais pour un couple lisse $(Y,S)$ on a calculé explicitement la valeur
des $(R^{q}f_{\ast})F$ pour la topologie étale (\hyperref[XVI.cor3.4]{3.4}, \hyperref[XVI.cor3.6]{3.6}, 
\hyperref[XVI.thm3.7]{3.7}), et un
calcul analogue et facile, que nous laissons au lecteur, donne les
résultats analogues pour la topologie classique. On en déduit le
théorème pour $f$ par une comparaison directe.

\section{Le théorème de finitude pour les préschémas algébriques en
caractéristique zéro}\etiquette[5]{XVI.sec5}

\begin{enonce*}{Théorème 5.1}\etiquette[5.1]{XVI.thm5.1}
Soient $k$ un corps de caractéristique zéro, et $f:X\rightarrow S$ un
morphisme de type fini de schémas localement de type fini sur
$k$. Soit $F$ un faisceau constructible d'ensembles (\resp de groupes
finis, \resp abélien de torsion) sur $X$. Alors les $R^{q}f_{\ast}F$
sont également constructibles pour $q=0$ (\resp pour $q=0$, $1$, \resp
pour tout $q$).
\end{enonce*}

On met l'hypothèse que $k$ soit de caractéristique $0$ parce qu'on va
se servir du théorème de résolution des singularités [1]. La
démonstration vaut également pour la caractéristique $p>0$ si l'on
admet la résolution des singularités, \emph{pourvu que $F$ soit
d'ordres premiers à $p$}. On peut donc déduire un résultat en
caractéristique $\neq 0$ pour $\dim X\leq 2$, en appliquant les
résultats d'Abhyankar \cite{SGA4-16:3}. Un autre ingrédient de la démonstration, en
plus du sempiternel th.\pageoriginale de changement de base pour un
morphisme propre, est le « théorème de pureté cohomologique
absolu » \hyperref[XVI.cor3.9]{3.9}\nde{Travaux de Deligne et Gabber XXX}.

Rappelons que le théorème est déjà démontré pour un morphisme propre
et pour $S$ arbitraire (\hyperrefext[XIV][XIV.thm1.1]{1.1}). D'ailleurs on
obtiendra dans Exposé XIX % syntaxe, étiquette
(encore en car. $0$) le théorème de finitude sous les hypothèses
beaucoup plus faibles que $S$ soit un schéma \emph{excellent} et que
$f$ soit de type fini, en prouvant pour de tels schémas le théorème de
pureté absolu sous la forme signalée dans \hyperref[XVI.rems3.10]{3.10}.

Dans le cas de caractéristique $\neq 0$ et d'un schéma de dimension
$\geq 3$, on n'a pour l'instant que la conséquence facile suivante de
\hyperrefext[XI][XI.prop3.3]{3.3}:

\begin{enonce*}{Théorème 5.2}\etiquette[5.2]{XVI.thm5.2}
Soit $X$ un schéma algébrique lisse sur $\Spec k$, $k$ un corps
séparablement clos, et soit $F$ un faisceau de groupes abéliens finis
(\resp de groupes finis) qui est \emph{localement constant,} et
d'ordres premiers à la caractéristique de $k$. Alors $H^{q}(X,F)$ est
un groupe fini (\resp un ensemble pointé fini) pour tout $q$ (\resp
pour $q=0,1$).
\end{enonce*}

\noindent
\emph{Démonstration de} \hyperref[XVI.thm5.2]{5.2}. On peut supposer $k$
algébriquement clos (\hyperrefext[VIII][VIII.thm1.1]{1.1}).
Récurrence sur $n=\dim X$: D'après \hyperrefext[XI][XI.prop3.3]{3.3} et il existe un
hyper-recouvrement de $X$ par des ouverts qui sont des « bons
voisinages ». On se réduit immédiatement, par la suite spectrale
d'un hyper-recouvrement, au cas où $X$ est un bon voisinage, donc
admet une fibration élémentaire (\hyperrefext[XI][XI.def3.1]{3.1})
$$
\xymatrix@=1.5cm{
X \ar[rr]^{j} \ar[rd]_{f}&         &
\bar{X}\ar[ld]^{\bar{f}} \\
             &      S  & .
}
$$
D'après\pageoriginale (\hyperref[XVI.prop3.2]{3.2}, \hyperref[XVI.cor3.6]{3.6}, \hyperref[XVI.thm3.7]{3.7}),
le faisceau $j_{\ast}F$ est
localement constant sur $X$, $(R^{1}j_{\ast})F$ est localement
constant sur $Y=\bar{X}-X$, et $(R^{q}j_{\ast})F=0$ si $q>0$. Par
suite il résulte de \hyperref[XVI.thm2.1]{2.1} que $(R^{p}\bar{f}_{\ast})(R^{q}j_{\ast})F$
est un faisceau localement constant sur $S$ pour les valeurs de $q$
envisagées, et le résultat résulte de l'hypothèse de récurrence et de
la suite spectrale de Leray (\resp de la suite exacte \hyperrefext[XII][XII.prop3.2]{3.2}).

La démonstration de \hyperref[XVI.thm5.1]{5.1} est très voisine 
de celle de \hyperref[XVI.thm4.1]{4.1} pour le cas
$F$ constructible. Nous nous bornerons à indiquer les grandes lignes:
\sisi{Récurrence}{récurrence} sur $\dim X=n$. On se réduit d'abord au cas $X$ et $S$
affines, donc $f$ quasi-projectif, en appliquant la méthode de la
démonstration de 1.3. Ensuite on applique \sisi{IX 2.15}{\hyperrefext[IX][IX.prop2.14]{2.14}} et (\hyperrefext[VIII][VIII.prop5.5]{5.5} et \hyperrefext[VIII][VIII.rem5.8]{5.8})
pour se ramener au cas $F$ constant. On a un diagramme commutatif
$$
\xymatrix@=1.5cm{
X \ar[rr]^{h} \ar[rd]_{f}&         &                  \bar{X}\ar[ld]^{\bar{f}}\\
             &      S  & ,                  
}
$$
où $h$ est une immersion ouverte dense et où $\bar{f}$ est propre,
donc le théorème est vrai pour $\bar{f}$ (\hyperrefext[XIV][XIV.thm1.1]{1.1}). Il suffit ainsi de
démontrer le théorème pour $(h,F)$, c'est-à-dire pour une immersion
ouverte dense\pageoriginale et un faisceau constant. Soit
$h:S'\rightarrow S$ une résolution des singularités de $S$ et
$$
\xymatrix@C=1.75cm@R=1cm{
X  \ar[d]^{f} & X' \ar[l]_{h'} \ar[d]^{f'}\\
S          & S' \ar[l]_{h} 
}
$$
le diagramme cartésien déduit de $h$. On se ramène à démontrer le
théorème pour $f'$, donc au cas où de plus $S=S'$ est régulier. Par la
méthode de récurrence employée dans la fin de la démonstration de \hyperref[XVI.thm4.1]{4.1},
on se réduit au cas où $Y=S-X$ est non-singulier, \ie où $(Y,S)$ est
un couple lisse, et on termine en appliquant (\hyperref[XVI.prop3.2]{3.2}, \hyperref[XVI.cor3.4]{3.4},
\hyperref[XVI.cor3.6]{3.6}, \hyperref[XVI.thm3.7]{3.7}).

\begin{enonce*}[remark]{Remarque 5.3}\etiquette[5.3]{XVI.rem5.3}
La démonstration de \hyperref[XVI.thm5.1]{5.1} qu'on vient d'esquisser est également valable
lorsqu'on se donne un anneau noethérien $A$ à gauche, annulé par un
entier $n>0$ (qu'il faut supposer premier à la caractéristique si
celle-ci n'est pas supposée nulle, \cf remarques suivant \hyperref[XVI.thm5.1]{5.1}), et
qu'on considère des faisceaux de $A$-modules à gauche constructibles.
\end{enonce*}

\begin{thebibliography}{99}
\bibitem{SGA4-16:1} Hironaka, H. Resolution of singularities of an
algebraic variety over a field of characteristic zero, Annals of
Math., vol. 79 (1964), p. 109.

\bibitem{SGA4-16:2} Godement,\pageoriginale P. Théorie des
\sisi{Faisceaux}{faisceaux}, Paris, 1958\sisi{}{.}

\bibitem{SGA4-16:3} Abhyankar, S. Local uniformization on algebraic
surfaces over ground fields of characteristic $p\neq 0$, Ann. of
Math., vol. 63 (1956), p. 491.

\bibitem{SGA4-16:4} J.P. Serre, \sisi{Géométrie Algébrique et Géométrie
Analytique}{Géométrie algébrique et géométrie analytique},
Annales Inst. Fourier 1956, p. 1-42 (\sisi{Cité}{cité} GAGA).
\end{thebibliography}

%XII 4.1 no ref. 78 
%XII 4.1.2 no ref. 992
%4.11.1 no ref. 1040
%XI 7.7 no ref. 1210
%IX 4.5 is a equation ref.
%IX 2.15 no ref. 1504

\ifx\danslelivre\undefined
\end{document}
\fi
